\subsection{NORMS AND TRACES RELATIVE TO A MODULE}

DEFINITION 1. Let M be an A-module admitting a finite basis as a K-module. For all \(a \in \mathbf{A}\) , let \(a_{\mathbf{M}}\) be the endomorphism \(x \mapsto ax\) of the K-module M. The trace, determinant and characteristic polynomial of \(a_{\mathbf{M}}\) are called respectively the trace, norm and characteristic polynomial of \(z\) relative to M.

The trace and norm of a are therefore elements of K, denoted respectively by \(\mathrm{Tr}_{\mathbf{M}/\mathbf{K}}(a)\) and \(\mathrm{N}_{\mathbf{M}/\mathbf{K}}(a)\) ; the characteristic polynomial of a is an element of K{[}X{]}, denoted by \(\mathrm{Pc}_{\mathbf{M}/\mathbf{K}}(a;\mathbf{X})\) . We omit K in the above notation when there is no risk of confusion.

From the properties of the trace and determinant of an endomorphism (II, § 4, no. 3 and § 8, no. 1) we obtain the relations

\[
\operatorname{Tr} _ {\mathbf {M}} (a + a ^ {\prime}) = \operatorname{Tr} _ {\mathbf {M}} (a) + \operatorname{Tr} _ {\mathbf {M}} (a ^ {\prime})\tag{1}
\]

\begin{enumerate}
\def\labelenumi{(\arabic{enumi})}
\setcounter{enumi}{1}
\tightlist
\item
\end{enumerate}

\[
\operatorname{Tr} _ {\mathbf {M}} (a a ^ {\prime}) = \operatorname{Tr} _ {\mathbf {M}} (a ^ {\prime} a)\tag{3}
\]

\[
\mathrm{N} _ {\mathbf {M}} (a a ^ {\prime}) = \mathrm{N} _ {\mathbf {M}} (a) \mathrm{N} _ {\mathbf {M}} (a ^ {\prime})
\]

for all \(a, a'\) in A.

Let \((e_i)_{1\leq i\leq n}\) be a basis of the K-module M and \((m_{ij}(a))\) the matrix of the endomorphism \(a_{\mathbf{M}}\) with respect to this basis. The functions \(m_{ij}\) are linear forms on the K-module A and

\begin{enumerate}
\def\labelenumi{(\arabic{enumi})}
\setcounter{enumi}{3}
\tightlist
\item
\end{enumerate}

\[
\operatorname{Tr} _ {\mathbf {M}} (a) = \sum_ {i = 1} ^ {n} m _ {i i} (a)\tag{5}
\]

\[
\mathrm{N} _ {\mathbf {M}} (a) = \det (m _ {i j} (a))\tag{6}
\]

\[
\operatorname{Pc} (a; \mathbf {X}) = \det (\delta_ {i j} \mathbf {X} - m _ {i j} (a)).
\]

It follows from the method of calculating a determinant (§ 8, no. 11, formula (50)) that

\[
\mathrm{Pc} _ {\mathbf {M}} (a; \mathbf {X}) = \mathbf {X} ^ {n} + c _ {1} \mathbf {X} ^ {n - 1} + \dots + c _ {n}\tag{7}
\]

where

\[
c _ {1} = - \operatorname{Tr} _ {\mathbf {M}} (a), \quad c _ {n} = (- 1) ^ {n} \mathrm{N} _ {\mathbf {M}} (a).\tag{8}
\]

For \(\lambda \in \mathbf{K}\) ,

\[
\operatorname{Tr} _ {\mathbf {M}} (\lambda) = n. \lambda , \quad \mathrm{N} _ {\mathbf {M}} (\lambda) = \lambda^ {n}, \quad \operatorname{Pc} _ {\mathbf {M}} (\lambda ; \mathbf {X}) = (\mathbf {X} - \lambda) ^ {n}.\tag{9}
\]

Let \(\mathbf{K}'\) be a commutative K-algebra. Let \(\mathbf{M}' = \mathbf{K}' \otimes_{\mathbf{K}} \mathbf{M}\) and \(\mathbf{A}' = \mathbf{K}' \otimes_{\mathbf{K}} \mathbf{A}\) , so that \(\mathbf{M}'\) has an A'-module structure (§ 4, Example 2). As a \(\mathbf{K}'\) -module \(\mathbf{M}'\) admits the basis consisting of the \(1 \otimes e_i\) ( \(1 \leqslant i \leqslant n\) ) and the matrix of \(a_{\mathbf{M}}\) with respect to \((e_i)\) is equal to the matrix of \((1 \otimes a)_{\mathbf{M}'}\) with respect to \((e_i')\) . Then

\[
\operatorname{Tr} _ {\mathbf {M} ^ {\prime}} (1 \otimes a) = \operatorname{Tr} _ {\mathbf {M}} (a). 1, \quad \mathrm{N} _ {\mathbf {M} ^ {\prime}} (1 \otimes a) = \mathrm{N} _ {\mathbf {M}} (a). 1\tag{12}
\]

\[
\mathrm{Pc} _ {\mathbf {M} ^ {\prime}} (1 \otimes a; \mathbf {X}) = \mathrm{Pc} _ {\mathbf {M}} (a; \mathbf {X}). 1
\]

for all \(a \in \mathbf{A}\) , where 1 denotes the unit element of \(\mathbf{K}'\) . If in particular we take \(\mathbf{K}' = \mathbf{K}[\mathbf{X}]\) , then

\[
\mathrm{Pc} _ {\mathbf {M} / \mathbf {K}} (a; \mathbf {X}) = \mathrm{N} _ {\mathbf {M} [ \mathbf {X} ] / \mathbf {K} [ \mathbf {X} ]} (\mathbf {X} - a).\tag{13}
\]

\subsection{PROPERTIES OF NORMS AND TRACES RELATIVE TO A MODULE}

If \(\mathbf{M}\) and \(\mathbf{M}'\) are two isomorphic A-modules with finite bases over \(\mathbf{K}\) , then, for all \(a \in \mathbf{A}\) ,

\[
\operatorname{Tr} _ {\mathbf {M} ^ {\prime}} (a) = \operatorname{Tr} _ {\mathbf {M}} (a), \quad \mathrm{N} _ {\mathbf {M} ^ {\prime}} (a) = \mathrm{N} _ {\mathbf {M}} (a), \quad \operatorname{Pc} _ {\mathbf {M} ^ {\prime}} (a; \mathbf {X}) = \operatorname{Pc} _ {\mathbf {M}} (a; \mathbf {X})\tag{14}
\]

for if \(f\) is an isomorphism of \(\mathbf{M}\) onto \(\mathbf{M}'\) , the matrix of \(a_{\mathbf{M}}\) with respect to a basis \(\mathbf{B}\) of \(\mathbf{M}\) over \(\mathbf{K}\) is the same as the matrix of \(a_{\mathbf{M}'}\) with respect to the basis \(f(\mathbf{B})\) of \(\mathbf{M}'\) .

PROPOSITION 1. Let \(\mathbf{M} = \mathbf{M}_0 \supset \mathbf{M}_1 \supset \cdots \supset \mathbf{M}_r = \{0\}\) be a decreasing sequence of submodules of an A-module \(\mathbf{M}\) such that each of the K-modules \(\mathbf{P}_i = \mathbf{M}_{i-1} / \mathbf{M}_i\) ( \(1 \leqslant i \leqslant r\) ) admits a finite basis. Then the K-module \(\mathbf{M}\) admits a finite basis and

\[
\operatorname{Tr} _ {\mathbf {M}} (a) = \sum_ {i = 1} ^ {r} \operatorname{Tr} _ {\mathrm{P} _ {i}} (a), \quad \mathrm{N} _ {\mathbf {M}} (a) = \prod_ {i = 1} ^ {r} \mathrm{N} _ {\mathrm{P} _ {i}} (a)\tag{15}
\]

\[
\mathrm{Pc} _ {\mathbf {M}} (a; \mathbf {X}) = \prod_ {i = 1} ^ {r} \mathrm{Pc} _ {\mathbf {P} _ {i}} (a; \mathbf {X}).
\]

Let \(B_{i}^{\prime}\) be a basis of \(P_{i}\) over K; then a system of representatives \(B_{i}\) of \(B_{i}^{\prime}\) (mod. \(M_{i}\) ) is a basis of a supplementary submodule of the K-module \(M_{i}\) in the K-module \(M_{i-1}\) (II, § 1, no. 11, Proposition 21). The union B of the \(B_{i}\) ( \(1 \leqslant i \leqslant r\) ) is a basis of M over K. Let \(X_{ii}\) be the matrix of the endomorphism \(a_{P_{i}}\) with respect to the basis \(B_{i}^{\prime}\) . It is immediate that the matrix of \(a_{M}\) with respect to B is of the form

\[
\left( \begin{array}{c c c c c} X _ {r r} & X _ {r, r - 1} & \ldots & X _ {r, 1} \\ 0 & X _ {r - 1, r - 1} & \ldots & X _ {r - 1, 1} \\ \cdot & \cdot & \cdot & \cdot & \cdot \\ 0 & 0 & \ldots & X _ {1 1} \end{array} \right)
\]

and the proposition follows from formulae (4), (5) and (6) of no. 1 and formula (31) of § 8, no. 6.

PROPOSITION 2. Let A, A' be two K-algebras, M an A-module and M' an A'-module. Suppose that M and M' are free K-modules of respective dimensions n and n' and consider M ⊗K M' as an (A ⊗K A')-module (§ 4, no. 3, Example 2). Then, for a \ensuremath{\in} A and a' \ensuremath{\in} A',

\begin{enumerate}
\def\labelenumi{(\arabic{enumi})}
\setcounter{enumi}{15}
\tightlist
\item
\end{enumerate}

\[
\operatorname{Tr} _ {\mathbf {M} \otimes \mathbf {M}} (a \otimes a ^ {\prime}) = \operatorname{Tr} _ {\mathbf {M}} (a) \operatorname{Tr} _ {\mathbf {M} ^ {\prime}} (a ^ {\prime})\tag{17}
\]

\[
\mathrm{N} _ {\mathbf {M} \otimes \mathbf {M} ^ {\prime}} (a \otimes a ^ {\prime}) = (\mathrm{N} _ {\mathbf {M}} (a)) ^ {n ^ {\prime}} (\mathrm{N} _ {\mathbf {M} ^ {\prime}} (a ^ {\prime})) ^ {n}.
\]

Formula (16) follows from II, § 4, no. 4, formula (26) and formula (17) from § 8, no. 6, formula (33).

\subsection{NORM AND TRACE IN AN ALGEBRA}

DEFINITION 2. Let A be a K-algebra which is a finite-dimensional free K-module. For every element \(a \in \mathbf{A}\) , the trace (resp. norm,† resp. characteristic polynomial) of a relative to A and K is the trace (resp. determinant, resp. characteristic polynomial) of the endomorphism \(x \mapsto ax\) of the K-module A.

The trace, norm and characteristic polynomial of \(a \in A\) relative to A and K are denoted by \(\operatorname{Tr}_{\mathbf{A}/\mathbf{K}}(a)\) , \(\operatorname{N}_{\mathbf{A}/\mathbf{K}}(a)\) and \(\operatorname{Pc}_{\mathbf{A}/\mathbf{K}}(a; \mathbf{X})\) ; we omit K and even A from this notation when there is no risk of confusion. Note that the trace (resp. norm, characteristic polynomial) of \(a \in A\) is just the trace (resp. norm, characteristic polynomial) of a relative to the A-module \(A_{s}\) .

Suppose that A is the product \(A_{1} \times A_{2} \times \cdots \times A_{m}\) of a finite number finite-dimensional algebras over K which are free K-modules. Using the above remark and Proposition 1 of no. 2, we have, for every element

\[
a = (a _ {1}, \dots , a _ {m}) \in \mathrm{A},
\]

\[
\mathrm{Tr} _ {\mathbf {A} / \mathbb {K}} (a) = \sum_ {i = 1} ^ {m} \mathrm{Tr} _ {\mathbf {A} _ {i} / \mathbb {K}} (a _ {i}), \qquad \mathrm{N} _ {\mathbf {A} / \mathbb {K}} (a) = \prod_ {i = 1} ^ {m} \mathrm{N} _ {\mathbf {A} _ {i} / \mathbb {K}} (a _ {i})\tag{18}
\]

\[
\mathrm{Pc} _ {\mathbf {A} / \mathbf {K}} (a; \mathbf {X}) = \prod_ {i = 1} ^ {m} \mathrm{Pc} _ {\mathbf {A} _ {i} / \mathbf {K}} (a _ {i}; \mathbf {X}).
\]

Similarly, Proposition 2 of no. 2 shows that if A and A' are two algebras, which are free K-modules, of finite dimensions n, \(n'\) respectively over K, then, for \(a \in A\) , \(a' \in A'\) .

\begin{enumerate}
\def\labelenumi{(\arabic{enumi})}
\setcounter{enumi}{18}
\tightlist
\item
\end{enumerate}

\[
\operatorname{Tr} _ {\mathbf {A} \otimes \mathbf {A} ^ {\prime}} (a \otimes a ^ {\prime}) = \operatorname{Tr} _ {\mathbf {A}} (a) \operatorname{Tr} _ {\mathbf {A} ^ {\prime}} (a ^ {\prime})\tag{20}
\]

\[
\mathrm{N} _ {\mathbf {A} \otimes \mathbf {A} ^ {\prime}} (a \otimes a ^ {\prime}) = (\mathrm{N} _ {\mathbf {A}} (a)) ^ {n ^ {\prime}} (\mathrm{N} _ {\mathbf {A} ^ {\prime}} (a ^ {\prime})) ^ {n}.
\]

Finally let A be a finite-dimensional algebra over K which is a free K-module, h a homomorphism of K into a commutative ring \(K'\) and \(A' = A_{(K')}\) the \(K'\) -algebra derived from A by extending the scalars by means of h. It follows from formula (12) of no. 1 that, for all \(a \in A\) ,

\[
\begin{array}{r l} \mathrm{Tr} _ {\mathbf {A} ^ {\prime} / \mathbf {K} ^ {\prime}} (1 \otimes a) & = h (\mathrm{Tr} _ {\mathbf {A} / \mathbf {K}} (a)), \quad \mathrm{N} _ {\mathbf {A} ^ {\prime} / \mathbf {K} ^ {\prime}} (1 \otimes a) = h (\mathrm{N} _ {\mathbf {A} / \mathbf {K}} (a)) \\ \mathrm{Pc} _ {\mathbf {A} ^ {\prime} / \mathbf {K} ^ {\prime}} (1 \otimes a; \mathbf {X}) & = \tilde {h} (\mathrm{Pc} _ {\mathbf {A} / \mathbf {K}} (a; \mathbf {X})) \end{array}\tag{21}
\]

where \(\tilde{h}\) is the homomorphism \(\mathbf{K}[\mathbf{X}] \to \mathbf{K}'[\mathbf{X}]\) derived from \(h\) . In particular, for \(\mathbf{K}' = \mathbf{K}[\mathbf{X}]\) , we obtain, writing \(\mathrm{A}[\mathbf{X}] = \mathrm{A} \otimes_{\mathbf{K}} \mathbf{K}[\mathbf{X}]\) ,

\[
\mathrm{Pc} _ {\mathbf {A} / \mathbf {K}} (a; \mathbf {X}) = \mathrm{N} _ {\mathbf {A} [ \mathbf {X} ] / \mathbf {K} [ \mathbf {X} ]} (\mathbf {X} - a).\tag{22}
\]

More generally, if \(\mathbf{K}'\) is a commutative \(\mathbf{K}\) -algebra and \(\mathbf{A}' = \mathbf{A} \otimes_{\mathbf{K}} \mathbf{K}'\) , then, for all \(x \in \mathbf{A}'\) ,

\[
\mathrm{Pc} _ {\mathbf {A} / \mathbf {K}} (a; x) = \mathrm{N} _ {\mathbf {A} ^ {\prime} / \mathbf {K} ^ {\prime}} (x - a).
\]

Examples. (1) Let A be a quadratic algebra over K of type \((\alpha, \beta)\) and \((e_1, e_2)\) a basis of type \((\alpha, \beta)\) (§ 2, no. 3). For \(x = \xi e_1 + \eta e_2\) , \(\operatorname{Tr}_{A/K}(x) = 2\xi + \beta \eta\) and \(\mathrm{N}_{A/K}(x) = \xi^2 + \beta \xi \eta - \alpha \eta^2\) ; these functions are therefore identical with the Cayley trace and norm of \(x\) (§ 2, no. 24).

\begin{enumerate}
\def\labelenumi{(\arabic{enumi})}
\setcounter{enumi}{1}
\item
  Let A be a quaternion algebra over K. A direct calculation allows us to verify that \(\mathrm{Tr}_{\mathbf{A} / \mathbb{K}}(x) = 2\mathrm{T}(x)\) and \(\mathrm{N}_{\mathbf{A} / \mathbb{K}}(x) = (\mathrm{N}(x))^2\) , where \(\mathbf{T}\) and \(\mathbf{N}\) are the Cayley trace and norm (§ 2, no. 4).
\item
  Let \(\mathbf{A} = \mathbf{M}_n(\mathbf{K})\) and let the canonical basis \((E_{ij})\) of \(\mathbf{A}\) (II, § 10, no. 3) be arranged in lexicographic order. Then it is immediately seen that for every matrix \(\mathbf{X} = \sum_{i,j} \xi_{ij} E_{ij}\) the matrix (of order \(n^2\) ) of the endomorphism \(Y \mapsto XY\) is of the form
\end{enumerate}

\[
\left( \begin{array}{c c c c} X & 0 & \ldots & 0 \\ 0 & X & \ldots & 0 \\ . & . & . & . \\ 0 & 0 & \ldots & X \end{array} \right)
\]

whence \(\operatorname{Tr}_{\mathbf{A} / \mathbb{K}}(X) = n.\operatorname{Tr}(X)\) and \(\mathrm{N}_{\mathbf{A} / \mathbb{K}}(X) = (\det (X))^n\)

\subsection{PROPERTIES OF NORMS AND TRACES IN AN ALGEBRA}

PROPOSITION 3. Let A be a K-algebra admitting a finite basis. For an element \(a \in \mathbf{A}\) to be invertible, it is necessary and sufficient that its \(\mathbf{N}_{\mathbf{A} / \mathbf{K}}(a)\) be invertible in \(\mathbf{K}\) .

If \(a\) admits an inverse \(a'\) in \(\mathbf{A}\) , then

\[
\mathrm{N} _ {\mathbf {A} / \mathbf {K}} (a) \mathrm{N} _ {\mathbf {A} / \mathbf {K}} (a ^ {\prime}) = \mathrm{N} _ {\mathbf {A} / \mathbf {K}} (a a ^ {\prime}) = \mathrm{N} _ {\mathbf {A} / \mathbf {K}} (1) = 1
\]

by formula (3) of no. 1. Conversely, if \(\mathbf{N}_{\mathbf{A} / \mathbb{K}}(a)\) is invertible, the endomorphism \(h: x \mapsto ax\) is bijective (§ 8, no. 2, Theorem 1). Then there exists \(a' \in A\) such that \(aa' = 1\) ; then \(h(a'a - 1) = aa'a - a = (aa' - 1)a = 0\) , whence \(aa' = 1\) since \(h\) is injective. Hence \(a'\) is the inverse of \(a\) .

PROPOSITION 4. Let A be a K-algebra admitting a finite basis. For all \(a \in \mathbf{A}\) , \(\operatorname{Pc}_{\mathbf{A} / \mathbf{K}}(a; a) = 0\) .

This follows immediately from the Cayley-Hamilton theorem (§ 8, no. 11, Proposition 20).

PROPOSITION 5. Let A be a K-algebra and m a two-sided ideal of A. Suppose that \(A_{0} = A/m\) is a free K-module of finite dimension n, that there exists an integer r \textgreater{} 0 such that \(m^{r} = \{0\}\) and that \(m^{i-1}/m^{i}\) is a free \(A_{0}\) -module of finite dimension \(s_{i}\) for \(1 \leqslant i \leqslant r\) . Let \(s = s_{1} + \cdots + s_{r}\) and for all \(a \in A\) let \(a_{0}\) denote the class of a mod. m. Then A is a free K-module of dimension ns and, for all \(a \in A\) ,

\[
\operatorname{Tr} _ {\mathbf {A}} (a) = s. \operatorname{Tr} _ {\mathbf {A} _ {0}} (a _ {0}), \quad \mathrm{N} _ {\mathbf {A}} (a) = (\mathrm{N} _ {\mathbf {A} _ {0}} (a _ {0})) ^ {s}\tag{23}
\]

\[
\mathrm{Pc} _ {\mathbf {A}} (a; \mathbf {X}) = (\mathrm{Pc} _ {\mathbf {A} _ {0}} (a _ {0}; \mathbf {X})) ^ {s}.
\]

By virtue of II, § 1, no. 13, Proposition 25, \(\mathfrak{m}^{i - 1} / \mathfrak{m}^{i}\) is a free K-module of dimension \(ns_i\) . Hence Proposition 1 of no. 2 can be applied with \(\mathbf{P}_i = \mathfrak{m}^{i - 1} / \mathfrak{m}^{i}\) ;

this shows in the first place that A is a free K-module of dimension \(n(s_{1} + \cdots + s_{r}) = ns\) . Moreover, the hypothesis implies that the A-module \(P_{i}\) is isomorphic to a direct sum of \(s_{i}\) submodules isomorphic to the A-module \(A_{0}\) ; by Proposition 1 of no. 2, therefore \(\mathrm{N}_{\mathrm{P}_{i}}(a) = \mathrm{N}_{\mathrm{A}_{0}}(a)^{s_{i}}\) ; finally therefore

\[
\mathrm{N} _ {\mathbf {A}} (a) = \mathrm{N} _ {\mathbf {A} _ {0}} (a) ^ {s}.
\]

In this formula \(\mathrm{N}_{\mathbf{A}_{0}}(a)\) is defined by considering \(A_{0}\) as a left A-module and it is equal to the determinant of the K-linear mapping \(x \mapsto ax\) of \(A_{0}\) into itself; but, as \(ax = a_{0}x\) for \(x \in A_{0}\) , \(\mathrm{N}_{\mathbf{A}_{0}}(a) = \mathrm{N}_{\mathbf{A}_{0}}(a_{0})\) , which completes the proof of formula (23) for the norm. The two others are shown analogously.

PROPOSITION 6. Let A be a commutative K-algebra admitting a finite basis over K and V an A-module admitting a finite basis over A. Then V admits a finite basis over K and for every A-endomorphism u of V, if \(u_{\mathbf{K}}\) is the mapping u considered as a K-endomorphism of V,

\[
\operatorname{Tr} (u _ {\mathbb {K}}) = \operatorname{Tr} _ {\mathbf {A} / \mathbb {K}} (\operatorname{Tr} (u)), \quad \det (u _ {\mathbb {K}}) = \mathrm{N} _ {\mathbf {A} / \mathbb {K}} (\det (u))\tag{24}
\]

\[
\operatorname{Pc} (u _ {\mathbf {K}}; \mathbf {X}) = \mathrm{N} _ {\mathbf {A} [ \mathbf {X} ] / \mathbf {K} [ \mathbf {X} ]} (\operatorname{Pc} (u; \mathbf {X})).
\]

Let \((a_{i})_{1\leqslant i\leqslant m}\) be a basis of A over K and \((e_{j})_{1\leqslant j\leqslant n}\) a basis of V over A; then \((a_{i}e_{j})\) is a basis of V over K (II, § 1, no. 13, Proposition 25). On the other hand the third of formulae (24) can be deduced from the second applied to the endomorphism \(\mathbf{X} - \bar{\boldsymbol{u}}\) of the A{[}X{]}-module A{[}X{]} \(\otimes_{\mathbf{A}}\mathrm{V}\) (§ 8, no. 10). It will therefore suffice to show the first two formulae in (24). We shall first establish the following lemma:

Lemma 1. Let \(\mathbf{X}_{ij}\) ( \(1 \leqslant i \leqslant n\) , \(1 \leqslant j \leqslant n\) ) be \(n^2\) indeterminates, \(X\) the square matrix \((\mathbf{X}_{ij})\) of order \(n\) and \(\mathbf{D}(\mathbf{X}_{11}, \ldots, \mathbf{X}_{nn}) \in \mathbf{Z}[\mathbf{X}_{11}, \ldots, \mathbf{X}_{nn}]\) the determinant \(\det(X)\) . On the other hand let A be a commutative ring, \(\mathbf{M}_{ij}\) ( \(1 \leqslant i \leqslant n\) , \(1 \leqslant j \leqslant n\) ) \(n^2\) matrices of order \(m\) over A, which are pairwise permutable, and M the square matrix of order \(mn\) over A which can be expressed as a square matrix of matrices (II, § 10, no. 7)

\[
M = \left( \begin{array}{c c c c} M _ {1 1} & M _ {1 2} & \ldots & M _ {1 n} \\ M _ {2 1} & M _ {2 2} & \ldots & M _ {2 n} \\ \cdot & \cdot & \cdot & \cdot \\ M _ {n 1} & M _ {n 2} & \ldots & M _ {n n} \end{array} \right)
\]

Then the determinant of M is equal to the determinant of the square matrix

\[
\mathrm{D} (M _ {1 1}, \dots , M _ {n n})
\]

of order m.

We proceed by induction on \(n\) , the cases \(n = 0\) and \(n = 1\) being trivial.

Let \(Z\) be a new indeterminate and \(N_{ij}\) the matrix \(M_{ij} + \delta_{ij}ZI_m\) ( \(\delta_{ij}\) the Kronecker index). If \(D^{ij}(X_{11},\ldots,X_{nn})\) is the cofactor of \(X_{ij}\) in the matrix \(X\) (§ 8, no. 6), then

\[
\mathrm{X} _ {j l} \mathrm{D} ^ {k l} (\mathrm{X} _ {1 1}, \dots , \mathrm{X} _ {n n}) = \delta_ {j k} \mathrm{D} (\mathrm{X} _ {1 1}, \dots , \mathrm{X} _ {n n})\tag{25}
\]

(§ 8, no. 6, formula (28)). We write \(N_{ij}' = \mathbf{D}^{ij}(N_{11}, \ldots, N_{nn})\) , which is a square matrix of order \(m\) over \(\mathrm{A}[Z]\) and consider the product \(N.U\) , where

\[
U = \left( \begin{array}{c c c c} N _ {1 1} ^ {\prime} & 0 & \ldots & 0 \\ N _ {1 2} ^ {\prime} & I _ {m} & \ldots & 0 \\ \cdot & \cdot & \cdot & \cdot \\ N _ {1 n} ^ {\prime} & 0 & \ldots & I _ {m} \end{array} \right),
\]

\[
N = \left( \begin{array}{c c c c} N _ {1 1} & N _ {1 2} & \ldots & N _ {1 n} \\ N _ {2 1} & N _ {2 2} & \ldots & N _ {2 n} \\ \cdot & \cdot & \cdot & \cdot \\ N _ {n 1} & N _ {n 2} & \ldots & N _ {n n} \end{array} \right)
\]

Performing this product in blocks (II, § 10, no. 5) and using formulae (25), we obtain

\[
N. U = \left( \begin{array}{c c c c} P & N _ {1 2} & \ldots & N _ {1 n} \\ 0 & N _ {2 2} & \ldots & N _ {2 n} \\ \cdot & \cdot & \cdot & \cdot \\ 0 & N _ {n 2} & \ldots & N _ {n n} \end{array} \right)
\]

where we have written \(P = \mathrm{D}(N_{11},\ldots ,N_{nn})\) .Let

\[
Q = \left( \begin{array}{c c c c} N _ {2 2} & \ldots & N _ {2 n} \\ \cdot & \cdot & \cdot & \cdot \\ N _ {n 2} & \ldots & N _ {n n} \end{array} \right)
\]

which is a matrix of order \(m(n - 1)\) ; then (§8, no. 6, formula (31)) (det \(N\) ) ( \(\det U\) ) = ( \(\det P\) ) ( \(\det Q\) ) and \(\det U = \det N_{11}'\) . But by the induction hypothesis, \(\det Q = \det(\mathrm{D}^{11}(N_{11},\ldots ,N_{nn})) = \det N_{11}'\) and by virtue of the definition of the \(N_{ij}\) , clearly \(\det Q\) is a polynomial in A{[}Z{]} of degree \(m(n - 1)\) , whose term in \(Z^{m(n - 1)}\) has coefficient 1; it follows immediately that \(\det Q\) is not a divisor of zero in the graded algebra A{[}Z{]}. We therefore conclude that \(\det N = \det(\mathrm{D}(N_{11},\ldots ,N_{nn}))\) in A{[}Z{]}; if we substitute 0 for Z in these polynomials, then \(\det M = \det(\mathrm{D}(M_{11},\ldots ,M_{nn}))\) .

Having shown this lemma, the K-module V is the direct sum of the K-modules \(Ae_{j}\) ( \(1 \leqslant j \leqslant n\) ); we write \(u(e_{j}) = \sum_{k=1}^{n} c_{jk} e_{k}\) . For every element \(xe_{j} \in Ae_{j}\) , with \(x \in A\) , the component of \(u(xe_{j})\) in \(Ae_{k}\) is \(xc_{jk} e_{k}\) ; it follows that the matrix of \(u_{k}\) with respect to the basis \((a_{i} e_{j})\) of the K-module V can be expressed in the form of a square matrix of matrices \((M_{jk})\) , where \(M_{jk}\) is the matrix of the K-linear mapping \(xe_{j} \mapsto xc_{jk} e_{k}\) of \(Ae_{j}\) into \(Ae_{k}\) with respect to the bases \((a_{i} e_{j})_{1 \leqslant i \leqslant m}\) and \((a_{i} e_{k})_{1 \leqslant i \leqslant m}\) of these two K-modules (II, § 10, no. 5). If for all \(t \in A\) , \(M(t)\) denotes the matrix, with respect to the basis \((a_{i})_{1 \leqslant i \leqslant m}\) of A over K, of the endomorphism \(x \mapsto xt\) of A, then \(M_{jk} = M(c_{jk})\) ; as \(t \mapsto M(t)\) is a ring homomorphism the matrices \(M_{jk}\) are permutable with one another. Then

\[
\det u _ {\mathrm{K}} = \det (\mathrm{D} (M _ {1 1}, \dots , M _ {n n}))
\]

by Lemma 1. But as \(t \mapsto M(t)\) is a ring homomorphism, \(\mathrm{D}(M_{11},\ldots ,M_{nn})\) is the matrix of the K-endomorphism \(x \mapsto x.\det (c_{jk})\) of A with respect to basis \((a_t)\) ; by definition its determinant is therefore \(\mathrm{N}_{\mathbf{A} / \mathbf{K}}(\det (u))\) , which proves the second formula of (24). On the other hand,

\[
\operatorname{Tr} \left(u _ {\mathrm{K}}\right) = \sum_ {j = 1} ^ {n} \operatorname{Tr} \left(M _ {j j}\right) = \sum_ {j = 1} ^ {n} \operatorname{Tr} _ {\mathrm{A} / \mathrm{K}} \left(c _ {j j}\right) = \operatorname{Tr} _ {\mathrm{A} / \mathrm{K}} \left(\sum_ {j = 1} ^ {n} c _ {j j}\right) = \operatorname{Tr} _ {\mathrm{A} / \mathrm{K}} (\operatorname{Tr} (u))
\]

and the proof of Proposition 6 is complete.

COROLLARY. Let A be a commutative K-algebra admitting a finite basis over K and B an A-algebra admitting a finite basis over A. Then B admits a finite basis over K, and, for all \(b \in B\) (``transitivity formulae'')

\[
\operatorname{Tr} _ {\mathbf {B} / \mathbf {K}} (b) = \operatorname{Tr} _ {\mathbf {A} / \mathbf {K}} (\operatorname{Tr} _ {\mathbf {B} / \mathbf {A}} (b)), \quad \mathrm{N} _ {\mathbf {B} / \mathbf {K}} (b) = \mathrm{N} _ {\mathbf {A} / \mathbf {K}} (\mathrm{N} _ {\mathbf {B} / \mathbf {A}} (b))\tag{26}
\]

\[
\mathrm{Pc} _ {\mathbf {B} / \mathbf {K}} (b; \mathbf {X}) = \mathrm{N} _ {\mathbf {A} [ \mathbf {X} ] / \mathbf {K} [ \mathbf {X} ]} (\mathrm{Pc} _ {\mathbf {B} / \mathbf {A}} (b; \mathbf {X})).
\]

This follows immediately from Proposition 6, setting \(\mathbf{V} = \mathbf{B}\) and \(u(x) = bx\) .

Remark. Suppose that the homomorphism \(\lambda \mapsto \lambda, 1\) of K into A is injective and let K be identified with its image in A; suppose that A admits a finite basis \((e_i)_{1 \leqslant i \leqslant n}\) as a K-module. Let \(s\) be an automorphism of A such that \(s(\mathbf{K}) = \mathbf{K}\) . Let \(a\) be an element of A; then, by transporting the structure

\begin{enumerate}
\def\labelenumi{(\arabic{enumi})}
\setcounter{enumi}{26}
\tightlist
\item
\end{enumerate}

\[
\operatorname{Tr} _ {\mathbf {A} / \mathbf {K}} (s (a)) = s (\operatorname{Tr} _ {\mathbf {A} / \mathbf {K}} (a))\tag{28}
\]

\[
\mathrm{N} _ {\mathbf {A} / \mathbf {K}} (s (a)) = s (\mathrm{N} _ {\mathbf {A} / \mathbf {K}} (a)).
\]

*Consider also a derivation D of A (§ 10, no. 2) such that D(K) ⊂ K and write \(D (e _ {i}) = \sum_ {j = 1} ^ {n} e _ {j} \mu_ {j i}\) where \(\mu_ {j i} \in \mathbf {K}\); write

\[
a e _ {i} = \sum_ {j = 1} ^ {n} e _ {j} \lambda_ {j i} \quad \text { with } \quad \lambda_ {j i} \in \mathbf {K}.
\]

Then

\[
\mathrm{D} (a) e _ {i} + a \mathrm{D} (e _ {i}) = \mathrm{D} (a e _ {i}) = \sum_ {j = 1} ^ {n} (\mathrm{D} (e _ {j}) \lambda_ {j i} + e _ {j} \mathrm{D} (\lambda_ {j i})).
\]

It follows that

\[
\mathbf {D} (a) e _ {i} = \sum_ {j = 1} ^ {n} e _ {j} \nu_ {j i}
\]

with \(\nu_{ji} = \mathrm{D}(\lambda_{ji}) + \sum_{k=1}^{n} (\mu_{jk}\lambda_{ki} - \lambda_{jk}\mu_{ki})\) . As \(\sum_{i,k} (\mu_{ik}\lambda_{ki} - \lambda_{ik}\mu_{ki}) = 0\) , there-

fore \(\mathrm{Tr}_{\mathbf{A} / \mathbb{K}}(\mathbf{D}(a)) = \sum_{i = 1}^{n}\mathbf{D}(\lambda_{ii})\) , in other words

\[
\operatorname{Tr} _ {\mathbf {A} / \mathbf {K}} (\mathbf {D} (a)) = \mathbf {D} (\operatorname{Tr} _ {\mathbf {A} / \mathbf {K}} (a)). *\tag{29}
\]

\subsection{DISCRIMINANT OF AN ALGEBRA}

DEFINITION 3. Let A be a K-algebra admitting a finite basis of n elements. The discriminant of a sequence \((x_{1},\ldots,x_{n})\) of n elements of A, with respect to K, denoted by \(\mathrm{D}_{\mathrm{A}/\mathrm{K}}(x_{1},\ldots,x_{n})\) , is the discriminant of the square matrix

\[
\left(\operatorname{Tr} _ {\mathbf {A} / \mathbb {K}} \left(x _ {i} x _ {j}\right)\right) _ {1 \leqslant i \leqslant n, 1 \leqslant j \leqslant n}.
\]

Consider first a basis \((e_i)_{1\leqslant i\leqslant n}\) of A over K and write

\[
e _ {i} e _ {j} = \sum_ {k = 1} ^ {n} c _ {i j k} e _ {k} \quad \text { with } \quad c _ {i j k} \in \mathbf {K}.\tag{30}
\]

Then \(\mathrm{Tr}_{\mathbf{A} / \mathbf{K}}(e_i) = \sum_{s=1}^{n} c_{iss}\) , whence \(\mathrm{Tr}_{\mathbf{A} / \mathbf{K}}(e_i e_j) = \sum_{k,s} c_{ijk} c_{kss}\) and therefore

\[
\mathrm{D} _ {\mathrm{A} / \mathrm{K}} (e _ {1}, \dots , e _ {n}) = \det \left(\left(\sum_ {k, s} c _ {i j k} c _ {k s s}\right) _ {1 \leqslant i \leqslant n, 1 \leqslant j \leqslant n}\right).\tag{31}
\]

Now let \((x_{i})_{1\leqslant i\leqslant n},(x_{i}^{\prime})_{1\leqslant i\leqslant n}\) be two sequences of \(n\) elements of A and suppose that there exists a square matrix of order \(n\) , \(M = (m_{ij})\) , with coefficients in K, such that \(x_{i} = \sum_{j=1}^{n} m_{ij} x_{j}'\) for \(1 \leqslant i \leqslant n\) . We write

\[
T = \left(\operatorname{Tr} _ {\mathrm{A} / \mathrm{K}} \left(x _ {i} x _ {j}\right)\right) _ {1 \leqslant i \leqslant n, 1 \leqslant j \leqslant n}, \quad T ^ {\prime} = \left(\operatorname{Tr} _ {\mathrm{A} / \mathrm{K}} \left(x _ {i} ^ {\prime} x _ {j} ^ {\prime}\right)\right) _ {1 \leqslant i \leqslant n, 1 \leqslant j \leqslant n}.
\]

Then \(\mathrm{Tr}_{\mathbf{A} / \mathbf{K}}(x_i x_j) = \sum_{p,q} m_{ip} m_{jq} \mathrm{Tr}_{\mathbf{A} / \mathbf{K}}(x_p' x_q')\) , whence \(T = M. T'.^t M\) ; the rule of multiplication of determinants therefore gives

\[
\det T = \det M. \det T ^ {\prime}. \det ^ {t} M = (\det M) ^ {2} \det T ^ {\prime}
\]

whence finally

\[
\mathrm{D} _ {\mathrm{A} / \mathrm{K}} (x _ {1}, \dots , x _ {n}) = (\det M) ^ {2} \mathrm{D} _ {\mathrm{A} / \mathrm{K}} (x _ {i} ^ {\prime}, \dots , x _ {n} ^ {\prime}).\tag{32}
\]

The above formula shows in particular that the discriminants of two bases of A over K differ by the square of an invertible element of K and therefore generate the same (principal) ideal of K. This ideal \(\Delta_{A/K}\) is called the discriminant ideal of A over K; by formula (32) the discriminant of every sequence of n elements of A which differ only in the order of the terms have the same discriminant, for the determinant of a permutation matrix is equal to \(\pm1\) .

Examples. (1) If A is a quadratic algebra of type \((\alpha, \beta)\) over K, then (in the notation of §2, no. 3) \(\operatorname{Tr}(e_1) = 2\) , \(\operatorname{Tr}(e_2) = \beta\) ,

\[
\operatorname{Tr} (e _ {2} ^ {2}) = \alpha \operatorname{Tr} (e _ {1}) + \beta \operatorname{Tr} (e _ {2}) = 2 \alpha + \beta^ {2},
\]

whence \(\mathbf{D}_{\mathbf{A} / \mathbb{K}}(e_1,e_2) = \beta^2 +4\alpha .\)

\begin{enumerate}
\def\labelenumi{(\arabic{enumi})}
\setcounter{enumi}{1}
\item
  Let \(\mathbf{A} = \mathbf{K}[\mathbf{X}] / \mathbf{K}[\mathbf{X}]\mathbf{P}\) , where \(\mathrm{P}(\mathbf{X}) = \mathbf{X}^3 + p\mathbf{X} + q\) , so that if \(x\) is the image of \(\mathbf{X}\) in \(\mathbf{A}, 1, x, x^2\) form a basis of \(\mathbf{A}\) over \(\mathbf{K}\) and \(x^3 = -px - q\) . It is seen immediately that \(\operatorname{Tr}(1) = 3\) , \(\operatorname{Tr}(x) = 0\) , \(\operatorname{Tr}(x^2) = -2p\) , taking account of the relation \(x^3 = -px - q\) , \(\operatorname{Tr}(x^3) = -3q\) and \(\operatorname{Tr}(x^4) = 2p^2\) , whence easily \(\mathrm{D}_{\mathbf{A} / \mathbf{K}}(1, x, x^2) = -4p^3 - 27q^2\) .
\item
  Let A be a quaternion algebra of type \((\alpha, \beta, \gamma)\) over K and \((1, i, j, k)\) a basis of A of type \((\alpha, \beta, \gamma)\) ; taking account of § 3, no. 5, formula (30), it is easily found that \(\operatorname{Tr}(1) = 4\) , \(\operatorname{Tr}(i) = 2\beta\) , \(\operatorname{Tr}(j) = \operatorname{Tr}(k) = 0\) , then
\end{enumerate}

\[
\mathrm{D} _ {\mathrm{A/K}} (1, i, j, k) = - 16 \gamma^ {2} (\beta^ {2} + 4 \alpha) ^ {2}.
\]

\begin{enumerate}
\def\labelenumi{(\arabic{enumi})}
\setcounter{enumi}{3}
\tightlist
\item
  Let \(\mathbf{A} = \mathbf{M}_n(\mathbf{K})\) and consider the canonical basis \((E_{ij})_{1\leqslant i\leqslant n,1\leqslant j\leqslant n}\) of \(\mathbf{A}\) over \(\mathbf{K}\) (II, § 10, no. 3). It is immediate that \(\operatorname{Tr}_{\mathbf{A} / \mathbf{K}}(E_{ij}) = 0\) if \(j\neq i\) and \(\operatorname{Tr}_{\mathbf{A} / \mathbf{K}}(E_{ii}) = n\) for all \(i\) ; it follows without difficulty that the matrix \((\operatorname{Tr}(E_{ij}E_{hk}))\) of order \(n^2\) is of the form \(n.P\) , where \(P\) is a permutation matrix, whence \(\mathrm{D}_{\mathbf{A} / \mathbf{K}}((E_{ij})) = \pm n^{n^2}\) .
\end{enumerate}

\section{DERIVATIONS}

In this paragraph, and unless otherwise mentioned, the algebras considered are not assumed to be associative nor necessarily to possess a unit element; K denotes a commutative ring.

\subsection{COMMUTATION FACTORS}

When in this paragraph we speak of graduations without specifying them, we shall always mean graduations of type \(\Delta\) , where \(\Delta\) is a commutative group written additively. In this paragraph, a commutation factor over \(\Delta\) with values in Z is called a commutation factor over \(\Delta\) ( \(\S\) 4, no. 7, Definition 6). A commutation factor over \(\Delta\) is therefore identified with a mapping \(\varepsilon\colon(\alpha,\beta)\mapsto\varepsilon_{\alpha\beta}=\varepsilon(\alpha,\beta)\) of \(\Delta\times\Delta\) into the multiplicative group \(\{-1,1\}\) such that for \(\alpha,\alpha^{\prime},\beta,\beta^{\prime}\) in \(\Delta\) ,

\[
\left\{ \begin{array}{l} \varepsilon (\alpha + \alpha^ {\prime}, \beta) = \varepsilon (\alpha , \beta) \varepsilon (\alpha^ {\prime}, \beta) \\ \varepsilon (\alpha , \beta + \beta^ {\prime}) = \varepsilon (\alpha , \beta) \varepsilon (\alpha , \beta^ {\prime}) \\ \varepsilon (\beta , \alpha) = \varepsilon (\alpha , \beta). \end{array} \right.\tag{1}
\]

It follows that \(\varepsilon(2\alpha, \beta) = \varepsilon(\alpha, 2\beta) = 1\) .

When \(\Delta = \mathbf{Z}\) , every commutation factor \(\varepsilon\) is determined by giving \(\varepsilon(1, 1)\) ; there are therefore only two such factors, the first defined by

\[
\varepsilon (p, q) = 1 \quad \text { for } \quad p, q \text { in } \mathbf {Z}\tag{2}
\]

and the second by

\[
\varepsilon (p, q) = (- 1) ^ {p q} \quad \text { for } \quad p, q \text { in } \mathbf {Z}.\tag{3}
\]

\subsection{GENERAL DEFINITION OF DERIVATIONS}

Consider a commutative ring K, six graded K-modules of type \(\Delta\) : A, A', A'', B, B', B'', and three K-linear mappings

\[
\mu \colon \mathrm{A} \times \mathrm{A} ^ {\prime} \rightarrow \mathrm{A} ^ {\prime \prime}, \quad \lambda_ {1} \colon \mathrm{B} \times \mathrm{A} ^ {\prime} \rightarrow \mathrm{B} ^ {\prime \prime}, \quad \lambda_ {2} \colon \mathrm{A} \times \mathrm{B} ^ {\prime} \rightarrow \mathrm{B} ^ {\prime \prime}
\]

such that the corresponding K-linear mappings

\[
\mathrm{A} \otimes_ {\mathrm{K}} \mathrm{A} ^ {\prime} \rightarrow \mathrm{A} ^ {\prime \prime}, \quad \mathrm{B} \otimes_ {\mathrm{K}} \mathrm{A} ^ {\prime} \rightarrow \mathrm{B} ^ {\prime \prime}, \quad \mathrm{A} \otimes_ {\mathrm{K}} \mathrm{B} ^ {\prime} \rightarrow \mathrm{B} ^ {\prime \prime}
\]

are graded of degree 0. The image \(\mu(a, a')\) for \(a \in A\) , \(a' \in A'\) is simply denoted by \(a, a'\) or even \(aa'\) , and similarly for the two other bilinear mappings. The degree of \(a, a'\) is therefore the sum of the degrees of \(a\) and \(a'\) .

DEFINITION 1. Given the above and a commutation factor \(\varepsilon\) on \(\Delta \times \Delta\) , an \(\varepsilon\) -derivation (or \((\mathbf{K}, \varepsilon)\) -derivation) of degree \(\delta \in \Delta\) of \((\mathbf{A}, \mathbf{A}', \mathbf{A}'')\) into \((\mathbf{B}, \mathbf{B}', \mathbf{B}'')\) is a triple of graded K-linear mappings of degree \(\delta\) :

\[
d \colon \mathrm{A} \rightarrow \mathrm{B}, \quad d ^ {\prime} \colon \mathrm{A} ^ {\prime} \rightarrow \mathrm{B} ^ {\prime}, \quad d ^ {\prime \prime} \colon \mathrm{A} ^ {\prime \prime} \rightarrow \mathrm{B} ^ {\prime \prime}
\]

such that, for every homogeneous element \(a \in A\) and every element \(a' \in A'\)

\[
d ^ {\prime \prime} (a. a ^ {\prime}) = (d a). a ^ {\prime} + \varepsilon_ {\delta , \deg (a)} a. (d ^ {\prime} a ^ {\prime}).\tag{4}
\]

It obviously suffices by linearity to verify relation (4) when \(a\) and \(a'\) run through respective generating systems of A and A'.

It is often convenient to denote the three mappings \(d, d', d''\) by the same letter \(d\) (which can be justified by denoting equally by \(d\) the graded K-linear mapping of degree \(\delta\) )

\[
(a, a ^ {\prime}, a ^ {\prime \prime}) \mapsto (d a, d ^ {\prime} a ^ {\prime}, d ^ {\prime \prime} a ^ {\prime \prime})
\]

of \(\mathbf{A} \oplus \mathbf{A}' \oplus \mathbf{A}''\) into \(\mathbf{B} \oplus \mathbf{B}' \oplus \mathbf{B}''\) . Relation (4) can then be written more simply

\[
d (a. a ^ {\prime}) = (d a). a ^ {\prime} + \varepsilon_ {\delta . \deg (a)} a. (d a ^ {\prime}).\tag{5}
\]

The \(\varepsilon\) -derivations of (A, A', A'') into (B, B', B'') of given degree form a sub-K-module of the K-module of graded linear mappings

\[
\operatorname{Homgr} _ {\mathbb {K}} (\mathrm{A} \oplus \mathrm{A} ^ {\prime} \oplus \mathrm{A} ^ {\prime \prime}, \mathrm{B} \oplus \mathrm{B} ^ {\prime} \oplus \mathrm{B} ^ {\prime \prime}).
\]

When \(\varepsilon(\alpha, \beta) = 1\) for all \(\alpha, \beta\) in \(\Delta\) , we say simply derivation (or K-derivation) instead of \(\varepsilon\) -derivation. Derivations form a sub-K-module of

\[
\operatorname{Hom} _ {\mathbb {K}} (\mathrm{A} \oplus \mathrm{A} ^ {\prime} \oplus \mathrm{A} ^ {\prime \prime}, \mathrm{B} \oplus \mathrm{B} ^ {\prime} \oplus \mathrm{B} ^ {\prime \prime}).
\]

When \(\Delta = Z\) and \(\varepsilon(p, q) = (-1)^{pq}\) , every \(\varepsilon\) -derivation of even degree is a derivation; every \(\varepsilon\) -derivation of odd degree is often called an antiderivation (or K-antiderivation); an antiderivation d therefore satisfies

\[
d (a. a ^ {\prime}) = (d a). a ^ {\prime} + (- 1) ^ {\deg (a)} a. (d a ^ {\prime})\tag{6}
\]

for a homogeneous element \(a \in A\) .

Remarks. (1) The notion of derivation can be defined for non-graded modules by agreeing to give these modules the trivial graduation.

\begin{enumerate}
\def\labelenumi{(\arabic{enumi})}
\setcounter{enumi}{1}
\tightlist
\item
  If only \(\varepsilon\) -derivations of given degree \(\delta\) are considered, the commutation factor \(\varepsilon\) may be disposed of as follows: the bilinear mapping \(\lambda_2: \mathbf{A} \times \mathbf{B}' \to \mathbf{B}''\) is modified by replacing it by the bilinear mapping \(\lambda_2': \mathbf{A} \times \mathbf{B}' \to \mathbf{B}''\) such that, for every homogeneous \(a\) in \(\mathbf{A}\) and all \(b' \in \mathbf{B}'\) ,
\end{enumerate}

\[
\lambda_ {2} ^ {\prime} (a, b ^ {\prime}) = \varepsilon_ {\delta , \deg (a)} \lambda_ {2} (a, b ^ {\prime}).
\]

Then \(d\) is a derivation relative to the bilinear mappings \(\mu, \lambda_{1}, \lambda_{2}^{\prime}\) .

The general definition of \(\varepsilon\) -derivations given above is especially used in two cases:

Case (I): A = B, \(A' = B'\) , \(A'' = B''\) and the three bilinear mappings \(\mu\) , \(\lambda_{1}\) , \(\lambda_{2}\) are equal to the same mapping.

Case (II): \(\mathbf{A} = \mathbf{A}' = \mathbf{A}''\) , \(\mathbf{B} = \mathbf{B}' = \mathbf{B}''\) , so that (for \(\mu\) ) \(\mathbf{A}\) is a graded algebra and the two K-bilinear mappings.

\[
\lambda_ {1} \colon \mathrm{B} \times \mathrm{A} \rightarrow \mathrm{B}, \quad \lambda_ {2} \colon \mathrm{A} \times \mathrm{B} \rightarrow \mathrm{B}\tag{7}
\]

are such that the corresponding K-linear mappings \(B \otimes_{K} A \to B\) , \(A \otimes_{K} B \to B\) are graded of degree 0. An \(\varepsilon\) -derivation of degree \(\delta\) of A into B is then a graded K-linear mapping \(d: A \to B\) of degree \(\delta\) , such that for every homogeneous x in A and every \(y \in A\) , we have the relation

\[
d (x y) = (d x) y + \varepsilon_ {\delta , \deg (a)} x (d y).\tag{8}
\]

Consider in particular in case (II) the case where A is a unital associative K-algebra and \(\lambda_{1}\) and \(\lambda_{2}\) are the external laws of an (A, A)-bimodule (§ 4, no. 3, Definition 3). This holds notably when A and B are two unital associative K-algebras, a unital homomorphism of graded K-algebras \(\rho: A \to B\) is given and an (A, A)-bimodule structure is considered on B defined by the two external laws

\[
\lambda_ {1} \colon (b, a) \mapsto b \rho (a), \quad \lambda_ {2} \colon (a, b) \mapsto \rho (a) b
\]

for \(a\in \mathbf{A},b\in \mathbf{B}.\)

Cases (I) and (II) have the following case in common: consider a graded K-algebra A, take B = A, mappings (7) both being multiplication on A. We then speak of an ε-derivation (or (K, ε)-derivation) of the graded algebra A: it is a graded K-linear mapping of A into itself, of degree δ, satisfying (8) for every homogeneous x in A and all \(y \in A\) . In particular if A is a graded ring, considered as an (associative) Z-algebra, we speak of the ε-derivation of the ring A.

Let A be a unital commutative associative K-algebra and B an A-module; when we speak of a derivation of A into B, it will always be understood that we mean with the A-bimodule structure on B derived from its A-module structure; then the formula

\[
d (x y) = x (d y) + y (d x) \quad \text { for } \quad x \in \mathrm{A}, y \in \mathrm{A}\tag{9}
\]

holds for such a derivation \(d: A \rightarrow B\) .

\subsection{EXAMPLES OF DERIVATIONS}

Example 1. *Let A be an R-algebra of differentiable mappings of R into R and let \(x_0\) be a point of R; R can be considered as an A-module with the external law \((f, a) \mapsto f(x_0)a\) . Then the mapping \(f \mapsto \mathrm{D}f(x_0)\) is a derivation, since (Functions of a Real Variable, I, § 1, no. 3)

\[
(\mathbf {D} (f g)) (x _ {0}) = (\mathbf {D} f (x _ {0})) g (x _ {0}) + f (x _ {0}) (\mathbf {D} g (x _ {0})). *
\]

Example 2. *Let X be a differentiable manifold of class C∞ and let A be the graded R-algebra of differential forms on X. The mapping which associates with every differential form ω on X its exterior differential dω is an anti-derivation of degree +1 (Differentiable and Analytic Manifolds, R, § 8).*

Example 3. Let A be an associative K-algebra. For all \(a \in \mathbf{A}\) , the mapping \(x \mapsto ax - xa\) is a derivation of the algebra A (cf.~no. 6).

Example 4. Let M be a K-module and A the exterior algebra \(\bigwedge (\mathbf{M}^{*})\) with its usual graduation (§ 7, no. 1). *It will be seen in § 11, no. 9 that, for all \(x\in \mathbf{M}\) , the right interior product \(i(x)\) is an antiderivation of A of degree \(-1\).*

Example 5. Returning to the general situation of Definition 1 of no. 2, let \(\overline{\mathbf{K}}\) be another commutative ring and \(\rho: \mathbf{K} \to \overline{\mathbf{K}}\) a ring homomorphism; let \(\overline{\mathbf{A}}, \overline{\mathbf{A}'}\) , \(\overline{\mathbf{A}''}\) , \(\overline{\mathbf{B}}, \overline{\mathbf{B}'}\) , \(\overline{\mathbf{B}''}\) denote the graded \(\overline{\mathbf{K}}\) -modules obtained respectively from \(\mathbf{A}, \mathbf{A}', \mathbf{A}''\) , \(\mathbf{B}, \mathbf{B}'\) , \(\mathbf{B}''\) by extending the ring of scalars to \(\overline{\mathbf{K}}\) (II, § 11, no. 5); we derive from \(\mu\) , \(\lambda_1\) and \(\lambda_2\) \(\overline{\mathbf{K}}\) -bilinear mappings

\[
\bar {\mu}: \overline {{{A}}} \times \overline {{{A}}} ^ {\prime} \rightarrow \overline {{{A}}} ^ {\prime \prime}, \quad \bar {\lambda} _ {1}: \bar {B} \times \overline {{{A}}} ^ {\prime} \rightarrow \bar {B} ^ {\prime \prime}, \quad \bar {\lambda} _ {2}: \overline {{{A}}} \times \bar {B} ^ {\prime} \rightarrow \bar {B} ^ {\prime \prime}
\]

by considering the tensor products by \(l_{\bar{K}}\) of the corresponding K-linear mappings to \(\mu\) , \(\lambda_{1}\) and \(\lambda_{2}\) (II, § 5, no. 1). Then, if d is an \(\varepsilon\) -derivation of degree \(\delta\) of \((\mathbf{A}, \mathbf{A}', \mathbf{A}^{\prime\prime})\) into \((\mathbf{B}, \mathbf{B}', \mathbf{B}^{\prime\prime})\) , the mapping \(\bar{d} = d \otimes l_{\bar{K}}\) of \(\overline{A} \oplus \overline{A}' \oplus \overline{A}''\) into \(\overline{B} \oplus \overline{B}' \oplus \overline{B}''\) is an \(\varepsilon\) -derivation of degree \(\delta\) of \((\overline{A}, \overline{A}', \overline{A}^{\prime\prime})\) into \((\overline{B}, \overline{B}', \overline{B}'' )\) .

Example 6. Let A be a graded K-algebra of type Z; a graded linear K-mapping of degree 0, \(d: A \to A\) , is defined by taking, for \(x_{n} \in \mathbf{A}_{n}(n \in \mathbf{Z})\) , \(d(x_{n}) = nx_{n}\) . This mapping is a derivation of A, since, for \(x_{p} \in A_{p}, x_{q} \in A_{q}\) ,

\[
d (x _ {p} x _ {q}) = (p + q) x _ {p} x _ {q} = d (x _ {p}) x _ {q} + x _ {p} d (x _ {q}).
\]

\subsection{COMPOSITION OF DERIVATIONS}

We suppose in this no. that case (I) of no. 2 holds, that is that A, A', A'' are three graded K-modules of type Δ and that we are given a K-bilinear mapping μ: A × A' → A'' corresponding to a graded K-linear mapping of degree 0, A ⊗K A' → A''. The graded endomorphisms f of A ⊕ A' ⊕ A'' such that f(A) ⊂ A, f(A') ⊂ A' and f(A'') ⊂ A'' form a graded subalgebra of the graded associative algebra EndgrK(A ⊕ A' ⊕ A'') (§ 3, no. 1, Example 2). In particular two ε-derivations of (A, A', A'') can be composed, but it should not be thought that the composition of two ε-derivations is an ε-derivation.

On every graded algebra B of type \(\Delta\) is defined the \(\varepsilon\) -bracket (or simply bracket when \(\varepsilon = 1\) ) of two homogeneous elements \(u, v\) , by the formula (10)

\[
[ u, v ] _ {\varepsilon} = u v - \varepsilon_ {\deg u, \deg v} v u (\text { denoted   simply   by } [ u, v ] \text { if } \varepsilon = 1).
\]

By extending this mapping by linearity, a K-bilinear mapping \((u, v) \mapsto [u, v]_{\varepsilon}\) of B × B into B is obtained. Then, for homogeneous u and v in B

\[
[ v, u ] _ {\varepsilon} = - \varepsilon_ {\deg u, \deg v} [ u, v ] _ {\varepsilon}.
\]

Applying this definition to the graded algebra \(\operatorname{Endgr}_{\mathbf{K}}(\mathbf{A} \oplus \mathbf{A}' \oplus \mathbf{A}'')\) , the \(\varepsilon\) -bracket of two graded endomorphisms is thus defined.

PROPOSITION 1. Let \(d_1, d_2\) be two \(\varepsilon\) -derivations of \((\mathbf{A}, \mathbf{A}', \mathbf{A} ' )\) of respective degrees \(\delta_1, \delta_2\) . Then the \(\varepsilon\) -bracket

\[
[ d _ {1}, d _ {2} ] _ {\varepsilon} = d _ {1} \circ d _ {2} - \varepsilon_ {\delta_ {1}, \delta_ {2}} d _ {2} \circ d _ {1}
\]

is an \(\varepsilon\) -derivation of degree \(\delta_1 + \delta_2\) . Moreover, if \(d\) is an \(\varepsilon\) -derivation of \((\mathbf{A}, \mathbf{A}', \mathbf{A}'')\) of degree \(\delta\) and if \(\varepsilon_{\delta, \delta} = -1\) , then \(d^2 = d \circ d\) is a derivation.

Suppose \(x \in A\) is homogeneous of degree \(\xi\) ; for all \(y \in A'\) ,

\[
\begin{array}{r l} d _ {1} (d _ {2} (x y)) & = ((d _ {1} d _ {2}) (x)) y + \varepsilon_ {\delta_ {1}, \delta_ {2} + \xi} (d _ {2} x) (d _ {1} y) \\ & \quad + \varepsilon_ {\delta_ {2}, \xi} (d _ {1} x) (d _ {2} y) + \varepsilon_ {\delta_ {1} + \delta_ {2}, \xi} x ((d _ {1} d _ {2}) (y)) \end{array}\tag{11}
\]

taking account of formulae (1) of no. 1. Exchanging the roles of \(d_1\) and \(d_2\) , we obtain, after simplifications again using (1) (no. 1),

\[
\begin{array}{r l} (d _ {1} d _ {2}) (x y) - \varepsilon_ {\delta_ {1}, \delta_ {2}} (d _ {2} d _ {1}) (x y) = & ((d _ {1} d _ {2}) (x)) y - \varepsilon_ {\delta_ {1}, \delta_ {2}} ((d _ {2} d _ {1}) (x)) y \\ & + \varepsilon_ {\delta_ {1} + \delta_ {2}, \xi} x ((d _ {1} d _ {2}) (y)) \\ & - \varepsilon_ {\delta_ {1}, \delta_ {2}} \varepsilon_ {\delta_ {1} + \delta_ {2}, \xi} x ((d _ {2} d _ {1}) (y)) \end{array}
\]

that is, writing \(d = [d_1, d_2]_{\varepsilon}\) and \(\delta = \delta_1 + \delta_2\) ,

\[
d (x y) = (d x) y + \varepsilon_ {\delta , \xi} x (d y)
\]

which proves that \(d\) is an \(\varepsilon\) -derivation.

On the other hand, if, in (11), we let \(d_1 = d_2 = d\) , \(\delta_1 = \delta_2 = \delta\) and \(\varepsilon_{\delta,\delta} = -1\) , we obtain, since then \(\varepsilon_{\delta,\delta + \xi} = -\varepsilon_{\delta,\xi}\) by (1),

\[
d ^ {2} (x y) = \left(d ^ {2} x\right) y + \varepsilon_ {2 5, \xi} x \left(d ^ {2} y\right)
\]

and as \(\varepsilon_{2\delta, \xi} = 1\) it is seen that \(d^2\) is a derivation.

COROLLARY. Suppose that \(\Delta = \mathbf{Z}\) . Then:

\begin{enumerate}
\def\labelenumi{(\roman{enumi})}
\item
  The square of an antiderivation is a derivation.
\item
  The bracket of two derivations is a derivation.
\item
  The bracket of an antiderivation and a derivation of even degree is an anti-derivation.
\item
  If \(d_1\) and \(d_2\) are antiderivations, \(d_1d_2 + d_2d_1\) is a derivation.
\end{enumerate}

Under the hypotheses of the beginning of this no., consider now a finite sequence \(\mathbf{D} = (d_i)_{1 \leqslant i \leqslant n}\) of pairwise permutable derivations of \((\mathbf{A}, \mathbf{A}', \mathbf{A}'')\) . For every polynomial \(\mathrm{P}(\mathbf{X}_1, \ldots, \mathbf{X}_n)\) in the algebra \(\mathbf{K}[\mathbf{X}_1, \ldots, \mathbf{X}_n]\) , the element \(\mathrm{P}(d_1, \ldots, d_n)\) of \(\operatorname{Endgr}_{\mathbf{K}}(\mathbf{A} \oplus \mathbf{A}' \oplus \mathbf{A}'')\) is then defined (§ 2, no. 9); its abbreviated notation is \(\mathrm{P}(\mathbf{D})\) .

PROPOSITION 2. With the above hypotheses and notation, consider \(2n\) indeterminates \(\mathrm{T}_1, \ldots, \mathrm{T}_n\) , \(\mathrm{T}_1', \ldots, \mathrm{T}_n'\) and for every polynomial \(\mathbf{F} \in \mathbf{K}[\mathbf{X}_1, \ldots, \mathbf{X}_n]\) write \(\mathbf{F}(\mathbf{T}) = \mathbf{F}(\mathbf{T}_1, \ldots, \mathbf{T}_n)\) , \(\mathbf{F}(\mathbf{T}') = \mathbf{F}(\mathbf{T}_1', \ldots, \mathbf{T}_n')\) and

\[
\mathrm{F} (\mathbf {T} + \mathbf {T} ^ {\prime}) = \mathrm{F} (\mathrm{T} _ {1} + \mathrm{T} _ {1} ^ {\prime}, \dots , \mathrm{T} _ {n} + \mathrm{T} _ {n} ^ {\prime}).
\]

Suppose that

\[
\mathbf {P} (\mathbf {T} + \mathbf {T} ^ {\prime}) = \sum_ {i} \mathbf {Q} _ {i} (\mathbf {T}) \mathbf {R} _ {i} (\mathbf {T} ^ {\prime})
\]

where the \(\mathbf{Q}_i\) and \(\mathbf{R}_i\) belong to \(\mathbf{K}[\mathbf{X}_1,\ldots ,\mathbf{X}_n]\) . Then, for all \(x\in \mathbf{A}\) and \(y\in \mathbf{A}'\) ,

\[
\mathbf {P} (\mathbf {D}) (x y) = \sum_ {i} \left(\mathbf {Q} _ {i} (\mathbf {D}) x\right) (\mathbf {R} _ {i} (\mathbf {D}) y).\tag{12}
\]

We introduce \(n\) other indeterminates \(T_1'', \ldots, T_n''\) and consider the polynomial algebra \(K[T_1, \ldots, T_n, T_1', \ldots, T_n', T_1'', \ldots, T_n''] = B\) ; on the other hand we consider the K-module \(M\) of bilinear mappings of \(A \times A'\) into \(A''\) ; a B-module structure is defined on \(M\) by writing, for every K-bilinear mapping \(f \in M\) and \(1 \leqslant i \leqslant n\) ,

\[
\left\{ \begin{array}{l} (\mathrm{T} _ {i} f) (a, a ^ {\prime}) = f (d _ {i} a, a ^ {\prime}) \\ (\mathrm{T} _ {i} ^ {\prime} f) (a, a ^ {\prime}) = f (a, d _ {i} a ^ {\prime}) \\ (\mathrm{T} _ {i} ^ {\prime \prime} f) (a, a ^ {\prime}) = d _ {i} (f (a, a ^ {\prime})). \end{array} \right.\tag{13}
\]

Since the \(d_{i}\) are permutable with one another, it is seen that, for every polynomial \(\mathbf{F} \in \mathbf{K}[\mathbf{X}_{1}, \ldots, \mathbf{X}_{n}]\) , \((\mathbf{F}(\mathbf{T})f)(a, a') = f(\mathbf{F}(\mathbf{D})a, a')\) ,

\[
(\mathbf {F} (\mathbf {T} ^ {\prime}) f) (a, a ^ {\prime}) = f (a, \mathbf {F} (\mathbf {D}) a ^ {\prime})
\]

and \((\mathbf{F}(\mathbf{T}^{\prime \prime})f) = \mathbf{F}(\mathbf{D})(f(a,a^{\prime}))\) . Hence, to prove (12) it suffices to show that

\[
(\mathbf {P} (\mathbf {T} ^ {\prime \prime}) - \sum_ {i} \mathbf {Q} _ {i} (\mathbf {T}) \mathbf {R} _ {i} (\mathbf {T} ^ {\prime})) \cdot \mu = 0\tag{14}
\]

or also \((\mathrm{P}(\mathbf{T}^{\prime\prime}) - \mathrm{P}(\mathbf{T} + \mathbf{T}^{\prime}))\) . \(\mu = 0\) in the B-module M. Now, the hypothesis that the \(d_{i}\) are derivations can also be expressed by saying that, for \(1 \leqslant i \leqslant n\) ,

\[
\left(\mathrm{T} _ {i} ^ {\prime \prime} - \mathrm{T} _ {i} - \mathrm{T} _ {i} ^ {\prime}\right). \mu = 0\tag{15}
\]

in the B-module M. By considering successively the polynomials

\[
\mathrm{P} \left(\mathrm{T} _ {1} ^ {\prime \prime}, \mathrm{T} _ {2} ^ {\prime \prime}, \dots , \mathrm{T} _ {n} ^ {\prime \prime}\right) - \mathrm{P} \left(\mathrm{T} _ {1} + \mathrm{T} _ {1} ^ {\prime}, \mathrm{T} _ {2} ^ {\prime \prime}, \dots , \mathrm{T} _ {n} ^ {\prime \prime}\right)
\]

\[
\mathrm{P} \left(\mathrm{T} _ {1} + \mathrm{T} _ {1} ^ {\prime}, \mathrm{T} _ {2} ^ {\prime \prime}, \dots , \mathrm{T} _ {n} ^ {\prime \prime}\right) - \mathrm{P} \left(\mathrm{T} _ {1} + \mathrm{T} _ {1} ^ {\prime}, \mathrm{T} _ {2} + \mathrm{T} _ {2} ^ {\prime}, \dots , \mathrm{T} _ {n} ^ {\prime \prime}\right)
\]

\[
\mathrm{P} \left(\mathrm{T} _ {1} + \mathrm{T} _ {1} ^ {\prime}, \dots , \mathrm{T} _ {n - 1} + \mathrm{T} _ {n - 1} ^ {\prime}, \mathrm{T} _ {n} ^ {\prime \prime}\right) - \mathrm{P} \left(\mathrm{T} _ {1} + \mathrm{T} _ {1} ^ {\prime}, \dots , \mathrm{T} _ {n - 1} + \mathrm{T} _ {n - 1} ^ {\prime}, \mathrm{T} _ {n} + \mathrm{T} _ {n} ^ {\prime}\right)
\]

it is seen that the difference \(\mathbf{P}(\mathbf{T}^{\prime \prime}) - \mathbf{P}(\mathbf{T} + \mathbf{T}^{\prime})\) can be written in the form

\[
\sum_ {i = 1} ^ {n} \left(\mathbf {T} _ {i} ^ {\prime \prime} - \mathbf {T} _ {i} - \mathbf {T} _ {i} ^ {\prime}\right) \mathrm{G} _ {i} (\mathbf {T}, \mathbf {T} ^ {\prime}, \mathbf {T} ^ {\prime \prime})
\]

where the \(G_{i}\) are elements of B. Relation (14) is therefore an immediate consequence of relations (15).

COROLLARY (Leibniz's formula). Let \(d_{i}\) ( \(1 \leqslant i \leqslant n\) ) be \(n\) derivations of \((\mathbf{A}, \mathbf{A}', \mathbf{A}'' )\) which are permutable with one another. For all \(\alpha = (\alpha_{1}, \ldots, \alpha_{n}) \in \mathbf{N}^{n}\) , we write

\[
d ^ {\alpha} = d _ {1} ^ {\alpha_ {1}} d _ {2} ^ {\alpha_ {2}} \dots d _ {n} ^ {\alpha_ {n}}.\tag{16}
\]

Then, for \(x \in \mathbf{A}\) and \(y \in \mathbf{A}'\) ,

\[
d ^ {\alpha} (x y) = \sum_ {\beta + \gamma = \alpha} ((\beta , \gamma)) d ^ {\beta} (x) d ^ {\gamma} (y)\tag{17}
\]

where we have written (in the notation introduced at the beginning of the chapter)

\[
((\beta , \gamma)) = (\beta + \gamma)! / (\beta ! \gamma!).\tag{18}
\]

This follows immediately from the multinomial formula (I, § 8, no. 2)

\[
(\mathbf {T} + \mathbf {T} ^ {\prime}) ^ {\alpha} = \sum_ {\beta + \gamma = \alpha} ((\beta , \gamma)) \mathbf {T} ^ {\beta} \mathbf {T} ^ {\prime \gamma}
\]

and Proposition 2.

\subsection{DERIVATIONS OF AN ALGEBRA A INTO AN A-MODULE}

We suppose in this no. that Case (II) of no. 2 holds. Then there is a graded K-algebra A and a graded K-module E and also two K-linear mappings of degree 0

\[
\mathbf {E} \otimes_ {\mathbf {K}} \mathbf {A} \rightarrow \mathbf {E}, \quad \mathbf {A} \otimes_ {\mathbf {K}} \mathbf {E} \rightarrow \mathbf {E}
\]

denoted by

\[
x \otimes a \mapsto x. a \quad \text { and } \quad a \otimes x \mapsto a. x \quad \text { for   } a \in A \text {   and   } x \in E.
\]

PROPOSITION 3. Let \(d:A\to E\) be an \(\varepsilon\) -derivation of degree \(\delta\) . Then \(\operatorname{Ker}(d)\) is a graded subalgebra of A; if A admits a unit element, it belongs to \(\operatorname{Ker}(d)\) .

Clearly \(\operatorname{Ker}(d)\) is a graded sub-K-module of A; further, relation (8) of no. 2 shows that, if x and y are two homogeneous elements belonging to \(\operatorname{Ker}(d)\) , then \(d(xy) = 0\) and hence \(xy \in \operatorname{Ker}(d)\) . Finally, if A admits a unit element 1 (of degree 0, cf.~§3, no. 1), relation (8) of no. 2, where x and y are replaced by 1, gives \(d(1) = d(1) + d(1)\) and hence \(d(1) = 0\) .

COROLLARY. Let \(d_1\) and \(d_2\) be two \(\varepsilon\) -derivations from \(A\) to \(E\) of the same degree \(\delta\) . If \(d_1\) and \(d_2\) take the same values at each element of a generating system of the algebra \(A\) , then \(d_1 = d_2\) .

\(d_{1} - d_{2}\) is an \(\varepsilon\) -derivation of degree \(\delta\) , hence \(\operatorname{Ker}(d_1 - d_2)\) is a subalgebra of A which contains a generating system of A and hence is equal to A.

PROPOSITION 4. Let \(d:A\to E\) be an \(\varepsilon\) -derivation of degree \(\delta\) . Suppose that A has a unit element 1 and let \(x\) be a homogeneous element of A with an inverse \(x^{-1}\) in A. Then

\[
d (x ^ {- 1}) = - \varepsilon_ {\delta , \deg (x)} x ^ {- 1} ((d x) x ^ {- 1}) = - \varepsilon_ {\delta , \deg (x)} (x ^ {- 1} (d x)) x ^ {- 1}.\tag{19}
\]

We have \(d(xx^{-1}) = d(1) = 0\) (Proposition 3), whence

\[
(d x) x ^ {- 1} + \varepsilon_ {\delta , \deg (x)} x (d (x ^ {- 1})) = 0
\]

which proves the first formula of (19). On the other hand, \(x^{-1}\) is homogeneous of degree \(-\deg(x)\) and \(\varepsilon_{\delta,\deg(x)} = \varepsilon_{\delta,-\deg(x)}\) by formulae (1) of no. 1; writing \(d(x^{-1}x) = 0\) , the second formula of (19) is obtained similarly.

PROPOSITION 5. Suppose that A is an integral domain and let L be its field of fractions. Every derivation of A into a vector space E over L (considered as an A-module) can be extended uniquely to a derivation of L into E.

Let \(d\) be a derivation of \(\mathbf{A}\) into \(\mathbf{E}\) and \(\bar{d}\) a derivation of \(\mathbf{L}\) into \(\mathbf{E}\) extending \(d\) ; then, for \(u \in \mathbf{A}, v \in \mathbf{A}, v \neq 0\) , of necessity, by virtue of (19),

\[
\bar {d} (u / v) = v ^ {- 1} d u - u v ^ {- 2} d v\tag{20}
\]

which proves the uniqueness of \(\bar{d}\) . Conversely, we show that \(\bar{d}\) can be defined by formula (20); it must first be verified that if \(u / v = u' / v'\) the value of the right hand side of (20) does not change when \(u\) is replaced by \(u'\) and \(v\) by \(v'\) . Now, \(uv' = vu'\) , hence \(v'(du) + u(dv') = v(du') + u'(dv)\) and therefore \(v'(du - uv^{-1}dv) = v(du' - u'v'^{-1}dv')\) , since \(uv'v^{-1} = u'\) and \(u'v'^{-1}v = u\) . Thus a mapping \(\bar{d}: L \to E\) has been defined which extends \(d\) ; it is immediately verified that it is K-linear and a derivation.

PROPOSITION 6. Suppose that A is a unital associative graded K-algebra and E a graded (A, A)-bimodule. If \(d: \mathbf{A} \to \mathbf{E}\) is an \(\varepsilon\) -derivation of degree \(\delta\) , then, for every finite sequence \((x_i)_{1 \leqslant i \leqslant n}\) of homogeneous elements of A, of respective degrees \(\xi_i\) ( \(1 \leqslant i \leqslant n\) ),

\[
d \left(x _ {1} x _ {2} \dots x _ {n}\right) = \sum_ {i = 1} ^ {n} \varepsilon_ {\delta, \xi_ {1} + \dots + \xi_ {i - 1}} x _ {1} \dots x _ {i - 1} \left(d x _ {i}\right) x _ {i + 1} \dots x _ {n}.\tag{21}
\]

Formula (21) is trivial for \(n = 0\) and is proved by induction on \(n\) , taking account of (4) (no. 2).

COROLLARY. Suppose that A is a unital commutative associative algebra and E an A-module. If \(d: \mathbf{A} \to \mathbf{E}\) is a derivation, then, for every integer \(n \geqslant 0\) ,

\[
d (x ^ {n}) = n x ^ {n - 1} (d x) \quad \text { for all } x \in \mathbf {A}.\tag{22}
\]

It suffices to give A the trivial graduation and apply (21) with all the \(x_{i}\) equal to \(x\) .

We return to the general case of an \(\varepsilon\) -derivation \(d: \mathbf{A} \to \mathbf{E}\) of degree \(\delta\) . Let \(Z_{\varepsilon}\) be the set of \(a \in \mathbf{A}\) such that for every homogeneous component \(a_{\alpha}\) of \(a\) of degree \(\alpha\) , for every homogeneous \(x\) in \(\mathbf{E}\) ,

\[
x a _ {\alpha} = \varepsilon_ {\alpha , \deg (x)} a _ {\alpha} x.\tag{23}
\]

If A is a unital associative graded algebra and E a graded (A, A)-bimodule it follows immediately from this definition that \(Z_{\varepsilon}\) is a graded subalgebra of A containing the unit element.

PROPOSITION 7. Suppose that A is a unital associative graded algebra and E a graded (A, A)-bimodule. Let \(d: \mathbf{A} \to \mathbf{E}\) be an \(\varepsilon\) -derivation of degree \(\delta\) and \(a\) a homogeneous element of \(Z_{\varepsilon}\) of degree \(\alpha\) . Then the mapping \(x \mapsto a(dx)\) is an \(\varepsilon\) -derivation of degree \(\delta + \alpha\) .

We write \(d'(x) = a(dx)\) ; for \(x\) homogeneous of degree \(\xi\) in \(A\) and \(y \in A\) , by virtue of (23) and (1) (no. 1),

\[
\begin{array}{r l} d ^ {\prime} (x y) & = a ((d x) y) + \varepsilon_ {\delta , \xi} a (x (d y)) = (a (d x)) y + \varepsilon_ {\delta + \alpha , \xi} (x a) (d y) \\ & = (d ^ {\prime} x) y + \varepsilon_ {\delta + \alpha , \xi} x (d ^ {\prime} y). \end{array}
\]

Proposition 7 says that the K-module of \(\varepsilon\) -derivations of A into E is a graded \(Z_{\varepsilon}\) -module of type \(\Delta\) .

\subsection{DERIVATIONS OF AN ALGEBRA}

Let A be a graded K-algebra; for every homogeneous element \(a \in \mathbf{A}\) , let \(\operatorname{ad}_{\varepsilon}(a)\) , or simply \(\operatorname{ad}(a)\) if no confusion can arise, denote the K-linear mapping of A into A

\[
x \mapsto [ a, x ] _ {\varepsilon}
\]

(no. 4, formula (10)) which is graded of degree deg a.

PROPOSITION 8. Let A be a graded K-algebra.

\begin{enumerate}
\def\labelenumi{(\roman{enumi})}
\tightlist
\item
  For every \(\varepsilon\) -derivation \(d: \mathbf{A} \to \mathbf{A}\) and every homogeneous element \(a\) of \(\mathbf{A}\) ,
\end{enumerate}

\[
[ d, \mathrm{ad} _ {\varepsilon} (a) ] _ {\varepsilon} = \mathrm{ad} _ {\varepsilon} (d a).\tag{24}
\]

\begin{enumerate}
\def\labelenumi{(\roman{enumi})}
\setcounter{enumi}{1}
\item
  If the algebra \(\mathbf{A}\) is associative, \(\operatorname{ad}_{\varepsilon}(a)\) is an \(\varepsilon\) -derivation of \(\mathbf{A}\) of degree \(\deg(a)\) .
\item
  Suppose that \(d\) is of degree \(\delta\) , let \(\alpha = \deg a\) and write \(f = [d, \operatorname{ad}_{\varepsilon}(a)]_{\varepsilon}\) . For every homogeneous element \(x \in A\) of degree \(\xi\) , we have, by (1) (no. 1),
\end{enumerate}

\[
\begin{array}{r l} f (x) & = d (a x - \varepsilon (\alpha , \xi) x a) - \varepsilon_ {\delta , \alpha} (a (d x) - \varepsilon_ {\alpha , \delta + \xi} (d x) a) \\ & = (d a) x + \varepsilon_ {\delta , \alpha} a (d x) - \varepsilon_ {\alpha , \xi} (d x) a - \varepsilon_ {\delta + \alpha , \xi} x (d a) \\ & \quad - \varepsilon_ {\delta , \alpha} a (d x) + \varepsilon_ {\alpha , \xi} (d x) a \\ & = (d a) x - \varepsilon_ {\delta + \alpha , \xi} x (d a) = [ d a, x ] _ {\varepsilon}. \end{array}
\]

\begin{enumerate}
\def\labelenumi{(\roman{enumi})}
\setcounter{enumi}{1}
\tightlist
\item
  For all \(x\) homogeneous of degree \(\xi\) and all \(y\) homogeneous of degree \(\eta\) in \(\mathbf{A}\) ,
\end{enumerate}

\[
\begin{array}{r l} \mathrm{ad} _ {\varepsilon} (a) (x y) & = a (x y) - \varepsilon_ {\alpha , \xi + \eta} (x y) a \\ & = (a x - \varepsilon_ {\alpha , \xi} x a) y + \varepsilon_ {\alpha , \xi} x (a y - \varepsilon_ {\alpha , \eta} y a) \\ & = \mathrm{ad} _ {\varepsilon} (a) (x). y + \varepsilon_ {\alpha , \xi} x. \mathrm{ad} _ {\varepsilon} (a) (y) \end{array}
\]

taking account of (1) and the associativity of A.

When A is associative, \(\operatorname{ad}_{\varepsilon}(a)\) is called the inner \(\varepsilon\) -derivation of A defined by a.

COROLLARY. Let A be an associative graded algebra. For two homogeneous elements, a, b of A,

\[
[ \mathrm{ad} _ {\varepsilon} (a), \mathrm{ad} _ {\varepsilon} (b) ] _ {\varepsilon} = \mathrm{ad} _ {\varepsilon} ([ a, b ] _ {\varepsilon}).\tag{25}
\]

It suffices to replace \(d\) by \(\operatorname{ad}_{\varepsilon}(a)\) and \(\operatorname{ad}_{\varepsilon}(a)\) by \(\operatorname{ad}_{\varepsilon}(b)\) in (24).

If \(\deg a = \alpha\) , \(\deg b = \beta\) , formula (25) is equivalent to the following relation for every homogeneous element \(c \in \mathbf{A}\) of degree \(\gamma\)

\[
\varepsilon_ {\alpha , \gamma} [ a, [ b, c ] _ {\varepsilon} ] _ {\varepsilon} + \varepsilon_ {\beta , \alpha} [ b, [ c, a ] _ {\varepsilon} ] _ {\varepsilon} + \varepsilon_ {\gamma , \beta} [ c, [ a, b ] _ {\varepsilon} ] _ {\varepsilon} = 0\tag{26}
\]

called the Jacobi identity.

\subsection{FUNCTORIAL PROPERTIES}

In this no., all algebras are assumed to be associative and unital and every algebra homomorphism is assumed to be unital.

PROPOSITION 9. Let A, B be two graded K-algebras, E an (A, A)-bimodule and F a graded (B, B)-bimodule; let \(\rho: \mathbf{A} \to \mathbf{B}\) be a graded algebra homomorphism and \(\theta: \mathbf{E} \to \mathbf{F}\) a graded A-homomorphism of A-bimodules (relative to \(\rho\) ), of degree 0. Then:

\begin{enumerate}
\def\labelenumi{(\roman{enumi})}
\item
  For every \(\varepsilon\) -derivation \(d':\mathbf{B}\to\mathbf{F}\) , \(d'\circ\rho:\mathbf{A}\to\rho^{*}(\mathbf{F})\) is an \(\varepsilon\) -derivation of the same degree.
\item
  For every \(\varepsilon\) -derivation \(d: \mathbf{A} \to \mathbf{E}\) , \(\theta \circ d: \mathbf{A} \to \rho^{*}(\mathbf{F})\) is an \(\varepsilon\) -derivation of the the same degree.
\end{enumerate}

The two assertions follow immediately from the relations

\[
\begin{array}{l} d ^ {\prime} (\rho (x y)) = d ^ {\prime} (\rho (x) \rho (y)) = d ^ {\prime} (\rho (x)) \rho (y) + \varepsilon_ {\delta ^ {\prime}, \xi} \rho (x) d ^ {\prime} (\rho (y)) \\ \theta (d (x y)) = \theta ((d x) y + \varepsilon_ {\delta, \xi} x (d y)) = \theta (d x) \rho (y) + \varepsilon_ {\delta, \xi} \rho (x) \theta (d y) \end{array}
\]

for \(x \in \mathbf{A}\) homogeneous of degree \(\xi\) and \(y \in \mathbf{A}\) , \(\delta\) and \(\delta'\) denoting the respective degrees of \(d\) and \(d'\) .

COROLLARY. Let S be a generating system of the algebra A. In order that \(d' \circ \rho = \theta \circ d\) , it is necessary and sufficient that \(d'(\rho(x)) = \theta(d(x))\) for all \(x \in S\) .

This is an immediate consequence of Proposition 9 and no. 5, Corollary to Proposition 3.

Under the conditions of Proposition 9, we know that B has (by means of \(\rho\) ) an (A, A)-bimodule structure (II, § 1, no. 14, Example 1).

PROPOSITION 10. Under the conditions of Proposition 9, for an \(\varepsilon\) -derivation \(d': \mathbf{B} \to \mathbf{F}\) to be A-linear for the left (resp. right) A-module structures on B and \(\rho_*(\mathbf{F})\) , it is necessary and sufficient that \(d'\) be zero on the subalgebra \(\rho(\mathbf{A})\) of B.

We perform the proof for left A-module structures. For \(a \in \mathbf{A}\) , \(b \in \mathbf{B}\) ,

\[
d ^ {\prime} (\rho (a) b) = d ^ {\prime} (\rho (a)) b + \rho (a) d ^ {\prime} b
\]

and hence if \(d' \circ \rho = 0\) , \(d'\) is linear for the left A-module structures on B and \(\rho^*(\mathbf{F})\) . Conversely, if this is so, in particular

\[
d ^ {\prime} (\rho (a)) = d ^ {\prime} (\rho (a). 1) = \rho (a) d ^ {\prime} (1) = 0
\]

(no. 5, Proposition 3).

In particular let \(D_{K}(B,F)\) denote the K-module of derivations of B into F (no. 2); those among these derivations which are A-linear, in other words those which are zero on \(\rho(A)\) , form a sub-K-module of \(D_{K}(B,F)\) , denoted by \(D_{A,\rho}(B,F)\) or simply \(D_{A}(B,F)\) (obviously \(D_{K}(B,F)=D_{K,\phi}(B,F)\) , where \(\phi:K\to B\) is the homomorphism defining the K-algebra structure on B).

Now let A, B, C be three graded K-algebras, \(\rho:A\to B\) , \(\sigma:B\to C\) two graded algebra homomorphisms and G a graded (C, C)-bimodule; if \(D_{A}(B,G)\) , \(D_{B}(C,G)\) and \(D_{A}(C,G)\) denote the respective K-modules \(D_{A,\rho}(B,\sigma_{*}(G))\) , \(D_{B,\sigma}(C,G)\) and \(D_{A,\sigma_{o}\rho}(C,G)\) , \(D_{B}(C,G)\) is clearly a sub-K-module of \(D_{A}(C,G)\) since \(\sigma(\rho(A))\subset\sigma(B)\) .

PROPOSITION 11. Under the above conditions, there is an exact sequence of K-homomorphisms

\[
0 \rightarrow \mathrm{D} _ {\mathrm{B}} (\mathrm{C}, \mathrm{G}) \xrightarrow {u} \mathrm{D} _ {\mathrm{A}} (\mathrm{C}, \mathrm{G}) \xrightarrow {v} \mathrm{D} _ {\mathrm{A}} (\mathrm{B}, \mathrm{G})\tag{27}
\]

where u is the canonical injection and v the homomorphism \(d \mapsto d \circ \sigma\) (Proposition 9). The kernel of v is the set of derivations \(d: C \to G\) such that \(d(\sigma(b)) = 0\) for all \(b \in B\) , which is precisely the image of u.

\subsection{RELATIONS BETWEEN DERIVATIONS AND ALGEBRA HOMOMORPHISMS}

We suppose again in this no. that Case (II) of no. 2 holds and the graded K-algebra A is not assumed to be associative. Given an element \(\delta \in \Delta\) , consider the graded K-module E( \(\delta\) ) (II, §11, no. 2) such that

\[
(\mathbf {E} (\delta)) _ {\mu} = \mathbf {E} _ {\mu + \delta}
\]

for all \(\mu \in \Delta\) . We define on the graded K-module \(\mathbf{A} \oplus \mathbf{E}(\delta)\) a graded K-algebra structure by setting, for every homogeneous element \(a \in \mathbf{A}\) and arbitrary elements \(a' \in \mathbf{A}, x, x'\) in \(\mathbf{E}(\delta)\)

\[
(a, x) \left(a ^ {\prime}, x ^ {\prime}\right) = \left(a a ^ {\prime}, x. a ^ {\prime} + \varepsilon_ {\delta, \deg (a)} a. x ^ {\prime}\right);\tag{28}
\]

the verification of the fact that this multiplication defines a graded ring structure is immediate.

The projection \(p:(a,x)\mapsto a\) is called the augmentation of the algebra \(\mathbf{A}\oplus\mathbf{E}(\delta)\) and is a graded algebra homomorphism. The graded K-linear mappings \(g:\mathbf{A}\to\mathbf{A}\oplus\mathbf{E}(\delta)\) of degree 0 such that the composition

\[
\mathrm{A} \xrightarrow {g} \mathrm{A} \oplus \mathrm{E} (\delta) \xrightarrow {p} \mathrm{A}
\]

is the identity \(l_{A}\) are the mappings of the form \(x \mapsto (x, f(x))\) , where \(f: A \to E\) is a graded K-linear mapping of degree \(\delta\) .

PROPOSITION 12. For a graded K-linear mapping \(f:A\to E\) of degree \(\delta\) to be an \(\varepsilon\) -derivation, it is necessary and sufficient that the mapping \(x\mapsto(x,f(x))\) of A into \(\mathbf{A}\oplus\mathbf{E}(\delta)\) be a graded K-algebra homomorphism.

Using the fact that for \(x\) homogeneous in \(\mathbf{A}\) and \(y \in \mathbf{A}\)

\[
(x y, f (x y)) = (x, f (x)). (y, f (y)),
\]

we obtain, taking account of (28), the equivalent relation

\[
f (x y) = f (x). y + \varepsilon_ {\delta , \deg (x)} x. f (y),
\]

whence the proposition.

PROPOSITION 13. For the algebra \(\mathbf{A} \oplus \mathbf{E}(\delta)\) to be associative and unital, it is necessary and sufficient that \(\mathbf{A}\) be associative and unital, and that the mappings \((a, x) \mapsto a.x\) and \((a, x) \mapsto x\) . \(a\) define on \(\mathbf{E}\) an \((\mathbf{A}, \mathbf{A})\) -bimodule structure; the unit element of \(\mathbf{A} \oplus \mathbf{E}(\delta)\) is then \((1, 0)\) .

If an element \((u,m)\in\mathbf{A}\oplus\mathbf{E}(\delta)\) is written as unit element of this algebra, it is immediately found that u must be the unit element of A; writing \((u,m).(0,x)=(0,x).(u,m)=(0,x)\) , we obtain u.x=x.u=x for \(x\in E\) and, writing \((u,m).(u,0)=(u,0).(u,m)=(u,0)\) , we obtain m=0. The fact that A is associative when \(\mathbf{A}\oplus\mathbf{E}(\delta)\) is follows from the fact that the augmentation is a surjective homomorphism. The condition \((x.a') .a''=x.(a'.a'')\) is then equivalent to \(((0,x)(a',0))(a'',0)=(0,x)((a',0)(a'',0))\) and similarly the condition \(a.(a'.x)=(a.a').x\) is equivalent to

\[
(a, 0) ((a ^ {\prime}, 0) (0, x)) = ((a, 0) (a ^ {\prime}, 0)) (0, x);
\]

finally the condition \(a \cdot (x \cdot a') = (a \cdot x) \cdot a'\) is equivalent to

\[
(a, 0) ((0, x) (a ^ {\prime}, 0)) = ((a, 0) (0, x)) (a ^ {\prime}, 0).
\]

\subsection{EXTENSION OF DERIVATIONS}

PROPOSITION 14. Let A be a commutative ring, M an A-module, B the A-algebra \(\mathsf{T}(\mathbf{M})\) (resp. \(\mathsf{S}(\mathbf{M})\) , resp. \(\bigwedge (\mathbf{M}))\) and E a (B, B)-bimodule. Let \(d_0: \mathbf{A} \to \mathbf{E}\) be a derivation of the ring A into the A-module E and \(d_{1}:M\to E\) an additive group homomorphism such that, for all \(a\in A\) and all \(x\in M\) ,

\[
d _ {1} (a x) = a d _ {1} (x) + d _ {0} (a). x,\tag{29}
\]

and further, when \(\mathbf{B} = \mathbf{S}(\mathbf{M})\) ,

\[
x. d _ {1} (y) + d _ {1} (x). y = y. d _ {1} (x) + d _ {1} (y). x\tag{30}
\]

for all \(x, y\) in \(\mathbf{M}\) , and, when \(\mathbf{B} = \bigwedge (\mathbf{M})\) ,

\[
x. d _ {1} (x) + d _ {1} (x). x = 0\tag{31}
\]

for all \(x \in \mathbf{M}\) . Then there exists one and only one derivation \(d\) of \(\mathbf{B}\) (considered as a \(\mathbf{Z}\) -algebra) into the \((\mathbf{B}, \mathbf{B})\) -bimodule \(\mathbf{E}\) such that \(d|_{\mathbf{A}} = d_0\) and \(d|M = d_1\) .

We take on the \(\mathbf{Z}\) -module \(\mathbf{B} \oplus \mathbf{E}\) the associative \(\mathbf{Z}\) -algebra structure defined by

\[
(b, t) (b ^ {\prime}, t ^ {\prime}) = (b b ^ {\prime}, b t ^ {\prime} + t b ^ {\prime})
\]

which has (1, 0) as unit element (no. 8, Proposition 13). Under the canonical injection \(t \mapsto (0, t)\) , E is identified with a two-sided ideal of B ⊕ E such that \(\mathrm{E}^2 = \{0\}\) . On the other hand, the mapping \(h_0 \colon \mathbf{B} \oplus \mathbf{E}\) defined by \(h_0(a) = (a, d_0(a))\) is a unital ring homomorphism (no. 8, Proposition 12); under this mapping, B ⊕ E then becomes an A-algebra. Moreover, if, for all \(x \in \mathbf{M}\) , we write \(h_1(x) = (x, d_1(x))\) , it follows from the definition of \(h_0\) and (29) that \(h_1(ax) = h_0(a)h_1(x)\) ; in other words \(h_1\) is an A-linear mapping of M into B ⊕ E. Then there exists one and only one A-algebra homomorphism, \(h \colon \mathbf{B} \to \mathbf{B} \oplus \mathbf{E}\) such that \(h \mid \mathbf{M} = h_1\) (and necessarily \(h \mid \mathbf{A} = h_0\) ): for, if \(\mathbf{B} = \mathsf{T}(\mathbf{M})\) , this follows from § 5, no. 1, Proposition 1; if \(\mathbf{B} = \mathbb{S}(\mathbf{M})\) , condition (30) shows that \(h(x)h(y) = h(y)h(x)\) for all \(x, y\) in M and the conclusion follows from § 6, no. 1, Proposition 2; finally if \(\mathbf{B} = \bigwedge (\mathbf{M})\) , condition (31) shows that \((h(x))^2 = 0\) for all \(x \in \mathbf{M}\) , since \(x \wedge x = 0\) and the conclusion follows from § 7, no. 1, Proposition 1. The homomorphism \(h\) is such that the composition \(p \circ h \colon \mathbf{B} \to \mathbf{B}\) with the augmentation \(\rho \colon \mathbf{B} \oplus \mathbf{E} \to \mathbf{B}\) is the identity \(l_B\) , for \(p \circ h\) and \(l_B\) coincide by definition for the elements of A and those of M and the set of these elements is a generating system of B. We can therefore write \(h(b) = (b, d(b))\) for all \(b \in B\) and the mapping \(b \mapsto d(b)\) of B into E is a derivation with the required properties, by virtue of Proposition 12 of no. 8.

COROLLARY. Let \(\mathbf{M}\) be a graded \(\mathbf{K}\) -module of type \(\Delta\) ; the \(\mathbf{K}\) -algebras \(\mathsf{T}(\mathbf{M})\) , \(\mathsf{S}(\mathbf{M})\) and \(\bigwedge (\mathbf{M})\) are given the corresponding graduations of type \(\Delta' = \Delta \times \mathbf{Z}\) (§ 5, no. 5, Proposition 7, § 6, no. 6, Proposition 10 and § 7, no. 7, Proposition 11). On the other hand \(\mathbf{M}\) is given the graduation of type \(\Delta'\) such that \(\mathbf{M}_{\alpha,1} = \mathbf{M}_{\alpha}\) for all \(\alpha \in \Delta\) and \(\mathbf{M}_{\alpha,n} = \{0\}\) for \(\alpha \in \Delta\) and \(n \neq 1\) . Let \(\varepsilon'\) be a commutation factor over \(\Delta'\) .

\begin{enumerate}
\def\labelenumi{(\roman{enumi})}
\item
  Let \(\mathbf{E}\) be a graded (left and right) \(\mathsf{T}(\mathbf{M})\) -bimodule of type \(\Delta'\) ; for all \(\delta \in \Delta\) and every integer \(n \in \mathbf{Z}\) , every graded \(\mathbf{K}\) -linear mapping \(f: \mathbf{M} \to \mathbf{E}\) of degree \(\delta_1' = (\delta, n)\) can be extended uniquely to an \(\varepsilon'\) -derivation \(d: \mathsf{T}(\mathbf{M}) \to \mathbf{E}\) of degree \(\delta'\) .
\item
  Let \(\mathbf{E}\) be a graded \(\mathbb{S}(\mathbf{M})\) -module of type \(\Delta'\) ; for a graded \(\mathbf{K}\) -linear mapping \(f: \mathbf{M} \to \mathbf{E}\) of degree \(\delta'\) to be extendable to an \(\varepsilon'\) -derivation \(d: \mathbb{S}(\mathbf{M}) \to \mathbf{E}\) of degree \(\delta'\) , it is necessary and sufficient that, for every ordered pair \((x, y)\) of homogeneous elements of \(\mathbf{M}\) ,
\end{enumerate}

\[
x. f (y) + \varepsilon_ {\delta^ {\prime}, (\deg (y), 1)} ^ {\prime} y. f (x) = y. f (x) + \varepsilon_ {\delta^ {\prime}, (\deg (x), 1)} ^ {\prime} x. f (y).\tag{33}
\]

The \(\varepsilon'\) -derivation \(d\) is then unique.

\begin{enumerate}
\def\labelenumi{(\roman{enumi})}
\setcounter{enumi}{2}
\tightlist
\item
  Let \(\mathbf{E}\) be a graded (left and right) \(\bigwedge (\mathbf{M})\) -bimodule of type \(\Delta'\) ; for a graded K-linear mapping \(f: \mathbf{M} \to \mathbf{E}\) of degree \(\delta'\) to be extendable to an \(\varepsilon'\) -derivation
\end{enumerate}

\[
d: \bigwedge (\mathbf {M}) \rightarrow \mathbf {E}
\]

of degree \(\delta'\) , it is necessary and sufficient that, for every homogeneous element \(x\) of \(\mathbf{M}\) ,

\[
x. f (x) + \varepsilon_ {\delta ^ {\prime}, (\deg (x), 1)} ^ {\prime} f (x). x = 0.\tag{34}
\]

The \(\varepsilon'\) -derivation \(d\) is then unique.

Remark 2 of no. 2 is applied with one of the external B-module laws on E (with B equal to T(M), S(M) or ∧(M)) modified; the external law thus modified is still, by (1) (no. 1), a B-module law and the B-module law thus obtained on E is still compatible with the other B-module structure. It then suffices to apply Proposition 14 with A = K and \(d_{0} = 0\) .

Example (1). In the application of Proposition 14 note that if \(d_0 = 0\) condition (29) means simply that \(d_1\) is A-linear. If we take in particular \(\mathbf{E} = \mathbf{B}\) and the (B, B)-bimodule structure derived from the ring structure on B, conditions (30) and (31) are automatically satisfied when \(d_1\) is taken to be the composition of an endomorphism \(s\) of \(\mathbf{M}\) and the canonical injection \(\mathbf{M} \to \mathbf{B}\) ; this is obvious for (30) since \(\mathbb{S}(\mathbf{M})\) is commutative and for (31) this follows from the fact that \(x\) and \(s(x)\) are of degree 1 in \(\bigwedge (\mathbf{M})\) . It is therefore seen that every endomorphism \(s\) of \(\mathbf{M}\) can be extended uniquely to a derivation \(\mathbf{D}_s\) of \(\mathbb{T}(\mathbf{M})\) (resp. \(\mathbb{S}(\mathbf{M})\) , resp. \(\bigwedge (\mathbf{M})\) ), which is of degree 0. Moreover, for two endomorphisms \(s, t\) of \(\mathbf{M}\) ,

\[
[ \mathbf {D} _ {s}, \mathbf {D} _ {t} ] = \mathbf {D} _ {[ s, t ]}\tag{35}
\]

for both sides are derivations of \(\mathsf{T}(\mathbf{M})\) (resp. \(\mathbb{S}(\mathbf{M})\) , resp. \(\bigwedge (\mathbf{M})\) ) which are equal to \([s, t]\) on \(\mathbf{M}\) .

The expression for \(D_{s}\) is obtained using formula (21) of no. 5, which gives respectively, for \(x_{1}, x_{2}, \ldots, x_{n}\) in M,

\[
\left\{ \begin{array}{l} \mathrm{D} _ {s} (x _ {1} \otimes x _ {2} \otimes \dots \otimes x _ {n}) \\ \qquad = \sum_ {i = 1} ^ {n} x _ {1} \otimes \dots \otimes x _ {i - 1} \otimes s (x _ {i}) \otimes x _ {i + 1} \otimes \dots \otimes x _ {n} \\ \mathrm{D} _ {s} (x _ {1} x _ {2} \dots x _ {n}) = \sum_ {i = 1} ^ {n} x _ {1} \dots x _ {i - 1} s (x _ {i}) x _ {i + 1} \dots x _ {n} \\ \mathrm{D} _ {s} (x _ {1} \wedge x _ {2} \wedge \dots \wedge x _ {n}) \\ \qquad = \sum_ {i = 1} ^ {n} x _ {1} \wedge \dots \wedge x _ {i - 1} \wedge s (x _ {i}) \wedge x _ {i + 1} \wedge \dots \wedge x _ {n}. \end{array} \right.\tag{36}
\]

In the case of the algebra \(\bigwedge (\mathbf{M})\) , there is the following interpretation of \(\mathbf{D}_s\) :

PROPOSITION 15. If M is a free K-module of finite rank n, then, for every endomorphism s of M, the restriction to \(\bigwedge^{n}(\mathbf{M})\) of the derivation \(\mathbf{D}_{s}\) is the homothety of ratio \(\operatorname{Tr}(s)\) .

Let \((e_j)_{1\leqslant j\leqslant n}\) be a basis of M and write \(e = e_1\wedge e_2\wedge \dots \wedge e_n\) . If

\[
s (e _ {j}) = \sum_ {k = 1} ^ {n} \alpha_ {j k} e _ {k},
\]

the third formula in (36) gives

\[
\mathrm{D} _ {s} (e) = \sum_ {i = 1} ^ {n} e _ {1} \wedge \dots \wedge e _ {i - 1} \wedge s \left(e _ {i}\right) \wedge e _ {i + 1} \wedge \dots \wedge e _ {n} = \left(\sum_ {j = 1} ^ {n} \alpha_ {j j}\right) e.
\]

Example (2). In the Corollary to Proposition 14, part (iii), let \(\Delta = \{0\}\) , the graduation on \(\bigwedge (\mathbf{M})\) then being the usual graduation of type \(\mathbf{Z}\) ; on the other hand take \(\varepsilon(p, q) = (-1)^{pq}\) . Then, for every linear form \(x^{*} \in \mathbf{M}^{*}\) on \(\mathbf{M}\) , \(x \mapsto \langle x, x^{*} \rangle\) is a graded K-linear mapping of degree \(-1\) of \(\mathbf{M}\) into \(\bigwedge (\mathbf{M})\) satisfying relation (34); then there exists an antiderivation \(i(x^{*})\) of \(\bigwedge (\mathbf{M})\) , of degree \(-1\) , such that (by virtue of formula (21) of no. 5)

\[
i \left(x ^ {*}\right) \left(x _ {1} \wedge \dots \wedge x _ {n}\right)
\]

\[
= \sum_ {i = 1} ^ {n} (- 1) ^ {i - 1} \left\langle x _ {i}, x ^ {*} \right\rangle x _ {1} \wedge \dots \wedge x _ {i - 1} \wedge x _ {i + 1} \wedge \dots \wedge x _ {n}
\]

and which is a special case of the inner product to be defined in § 11, no. 9, formula (68).

PROPOSITION 16. Let A be a commutative K-algebra, \(\mathbf{M}_i\) ( \(1 \leqslant i \leqslant n\) ) and P A-modules and H the A-module of A-multilinear mappings of \(\mathbf{M}_1 \times \mathbf{M}_2 \times \cdots \times \mathbf{M}_n\) into P. Suppose that there is given a K-derivation \(d_{0}:A\to A\) of the algebra A, for each i, a K-linear mapping \(d_{i}:M_{i}\to M_{i}\) and a K-linear mapping D:P→P, so that, for \(1\leqslant i\leqslant n\) , \((d_{0},d_{i},d_{i})\) is a K-derivation of \((\mathbf{A},\mathbf{M}_{i},\mathbf{M}_{i})\) into itself and \((d_{0},\mathbf{D},\mathbf{D})\) is a K-derivation of \((\mathbf{A},\mathbf{P},\mathbf{P})\) into itself. Then there exists a K-linear mapping \(D^{\prime}:H\to H\) such that \((d_{0},D^{\prime},D^{\prime})\) is a K-derivation of \((\mathbf{A},\mathbf{H},\mathbf{H})\) into itself and

\begin{enumerate}
\def\labelenumi{(\arabic{enumi})}
\setcounter{enumi}{36}
\tightlist
\item
  \(\mathbf{D}(f(x_1, \ldots, x_n))\)
\end{enumerate}

for all \(x_{i} \in \mathbf{M}_{i}\) for \(1 \leqslant i \leqslant n\) and \(f \in \mathbf{H}\) .

We show that for \(f \in H\) , the mapping \(D'f\) of \(M_{1} \times M_{2} \times \cdots \times M_{n}\) into P defined by (37) is A-multilinear. For \(a \in A\) ,

\[
\begin{array}{r l} (\mathrm{D} ^ {\prime} f) (a x _ {1}, x _ {2}, \ldots , x _ {n}) & = \mathrm{D} (a f (x _ {1}, \ldots , x _ {n})) - f (d _ {1} (a x _ {1}), x _ {2}, \ldots , x _ {n}) \\ & \quad - a \sum_ {i = 2} ^ {n} f (x _ {1}, \ldots , x _ {i - 1}, d _ {i} x _ {i}, x _ {i + 1}, \ldots , x _ {n}) \end{array}
\]

and by hypothesis

\[
\mathrm{D} (a f (x _ {1}, \dots , x _ {n})) = (d _ {0} a) f (x _ {1}, \dots , x _ {n}) + a \mathrm{D} (f (x _ {1}, \dots , x _ {n}))
\]

and \(d_{1}(ax_{1}) = (d_{0}a)x_{1} + a.d_{1}x_{1}\) , which gives

\[
(\mathrm{D} ^ {\prime} f) (a x _ {1}, x _ {2}, \dots , x _ {n}) = a. (\mathrm{D} ^ {\prime} f) (x _ {1}, \dots , x _ {n})
\]

and linearity in each of the \(x_{i}\) is proved similarly. On the other hand,

\[
\begin{array}{r l} (\mathrm{D} ^ {\prime} (a f)) (x _ {1}, \dots , x _ {n}) & = \mathrm{D} (a f (x _ {1}, \dots , x _ {n})) \\ & - \sum_ {i = 1} ^ {n} a f (x _ {1}, \dots , x _ {i - 1}, d _ {i} x _ {i}, x _ {i + 1}, \dots , x _ {n}) \\ & = (d _ {0} a) f (x _ {1}, \dots , x _ {n})) + a \mathrm{D} (f (x _ {1}, \dots , x _ {n})) \\ & - \sum_ {i = 1} ^ {n} a f (x _ {1}, \dots , x _ {i - 1}, d _ {i} x _ {i}, x _ {i + 1}, \dots , x _ {n}) \\ & = (d _ {0} a) f (x _ {1}, \dots , x _ {m}) + a (\mathrm{D} ^ {\prime} f) (x _ {1}, \dots , x _ {m}) \end{array}
\]

in other words

\[
\mathbf {D} ^ {\prime} (a f) = \left(d _ {0} a\right) f + a (\mathbf {D} ^ {\prime} f)
\]

which establishes the proposition.

Examples. (3) Applying Proposition 16 to the case \(n = 1\) , \(\mathbf{M}_1 = \mathbf{M}\) , \(P = A\) , then \(H = M^*\) , the dual of \(M\) , and it is seen that for a K-derivation \((d_0, d, d)\) of \((A, M, M)\) is derived a K-derivation \((d_0, d^*, d^*)\) of \((A, M^*, M^*)\) such that

\[
d _ {0} \langle m, m ^ {*} \rangle = \langle d m, m ^ {*} \rangle + \langle m, d ^ {*} m ^ {*} \rangle\tag{38}
\]

for \(m \in \mathbf{M}\) and \(m^* \in \mathbf{M}^*\) . The K-linear mapping of \(\mathbf{M} \oplus \mathbf{M}^*\) into itself which is equal to \(d\) on \(\mathbf{M}\) and to \(d^{*}\) on \(\mathbf{M}^{*}\) then satisfies condition (29) and there is therefore a K-derivation \(\mathbf{D}\) of the A-algebra \(\mathsf{T}(\mathbf{M} \oplus \mathbf{M}^{*})\) , which reduces to \(d_{0}\) on \(\mathbf{A}\) , to \(d\) on \(\mathbf{M}\) and to \(d^{*}\) on \(\mathbf{M}^{*}\) . The restriction \(d_{\mathbf{J}}^{\mathbf{I}}\) of \(\mathbf{D}\) to the sub-A-module \(\mathsf{T}_{\mathbf{J}}^{\mathbf{I}}(\mathbf{M})\) of \(\mathsf{T}(\mathbf{M} \oplus \mathbf{M}^{*})\) ( \(\S 5\) , no. 6) is a K-endomorphism of \(\mathsf{T}_{\mathbf{J}}^{\mathbf{I}}(\mathbf{M})\) such that \((d_{0}, d_{\mathbf{J}}^{\mathbf{I}}, d_{\mathbf{J}}^{\mathbf{I}})\) is a K-derivation of \((\mathbf{A}, \mathsf{T}_{\mathbf{J}}^{\mathbf{I}}(\mathbf{M}), \mathsf{T}_{\mathbf{J}}^{\mathbf{I}}(\mathbf{M}))\) . Moreover, for \(i \in \mathbf{I}\) , \(j \in \mathbf{J}\) , if we write \(\mathbf{I}' = \mathbf{I} - \{i\}\) , \(\mathbf{J}' = \mathbf{J} - \{j\}\) , it is immediately verified that for the contraction \(c_{j}^{t}\) ( \(\S 5\) , no. 6)

\[
c _ {j} ^ {t} (d _ {J} ^ {\mathrm{I}} (z)) = d _ {J ^ {\prime}} ^ {\mathrm{I} ^ {\prime}} (c _ {j} ^ {t} (z)) \quad \text { for   all } z \in \mathbf {T} _ {J} ^ {\mathrm{I}} (\mathbf {M}).
\]

\begin{enumerate}
\def\labelenumi{(\arabic{enumi})}
\setcounter{enumi}{3}
\tightlist
\item
  Let \(M_{i}\) ( \(1 \leqslant i \leqslant 3\) ) be three A-modules and, for each i, suppose that \((d_{0}, d_{i}, d_{i})\) is a derivation of \((\mathbf{A}, \mathbf{M}_{i}, \mathbf{M}_{i})\) ; applying Proposition 16 again for n = 1, a derivation \((d_{0}, d_{ij}, d_{ij})\) of \((\mathbf{A}, \mathbf{H}_{ij}, \mathbf{H}_{ij})\) , where \(\mathrm{H}_{ij} = \mathrm{Hom}_{\mathbf{A}}(\mathbf{M}_{i}, \mathbf{M}_{j})\) , is derived for each ordered pair \((i, j)\) . With this notation, for \(u \in \mathrm{Hom}_{\mathbf{A}}(\mathbf{M}_{1}, \mathbf{M}_{2})\) and \(v \in \mathrm{Hom}_{\mathbf{A}}(\mathbf{M}_{2}, \mathbf{M}_{3})\) ,
\end{enumerate}

\[
d _ {1 3} (v \circ u) = (d _ {2 3} v) \circ u + v \circ (d _ {1 2} u)\tag{39}
\]

as is immediately verified from the definitions.

\subsection{UNIVERSAL PROBLEM FOR DERIVATIONS; NON-COMMUTATIVE CASE}

Throughout the rest of §10 all the algebras are assumed to be associative and unital and all the algebra homomorphisms are assumed to be unital.

Let A be a K-algebra; the tensor product \(A \otimes_{K} A\) has canonically an (A, A)-bimodule structure under which

\[
x. (u \otimes v). y = (x u) \otimes (v y)\tag{40}
\]

for all \(x, y, u, v\) in A (§ 4, no. 3, Example 2). The K-linear mapping \(m: \mathrm{A} \otimes_{\mathbb{K}} \mathrm{A} \to \mathrm{A}\) corresponding to multiplication in A (and hence such that \(m(x \otimes y) = xy\) ) is an (A, A)-bimodule homomorphism; its kernel I is therefore a sub-bimodule of \(\mathrm{A} \otimes_{\mathbb{K}} \mathrm{A}\) .

Lemma 1. The mapping \(\delta_{\mathbf{A}}: x \mapsto x \otimes 1 - 1 \otimes x\) is a K-derivation of A into I and I is generated, as a left A-module, by the image of \(\delta_{\mathbf{A}}\) .

The first assertion follows from the fact that

\[
(x y) \otimes 1 - 1 \otimes (x y) = (x \otimes 1 - 1 \otimes x). y + x. (y \otimes 1 - 1 \otimes y)
\]

by (40). On the other hand, if the element \(\sum_{i} x_{i} \otimes y_{i}\) (for \(x_{i}, y_{i}\) in A) belongs to I, by definition \(\sum_{i} x_{i} y_{i} = 0\) and hence

\[
\sum_ {i} \left(x _ {i} \otimes y _ {i}\right) = \sum_ {i} x _ {i} (1 \otimes y _ {i} - y _ {i} \otimes 1)
\]

by (40), which completes the proof of the lemma.

PROPOSITION 17. The derivation \(\delta_{\mathbf{A}}\) has the following universal property: for every (A, A)-bimodule E and every K-derivation \(d: \mathbf{A} \to \mathbf{E}\) , there exists one and only one (A, A)-bimodule homomorphism \(f: \mathbf{I} \to \mathbf{E}\) such that \(d = f \circ \delta_{\mathbf{A}}\) .

Note first that, for every (A, A)-bimodule homomorphism \(f: I \to E, f \circ \delta_{A}\) is a derivation (no. 7, Proposition 9). Conversely, let \(d: A \to E\) be a K-derivation; then we prove first that if there exists an (A, A)-bimodule homomorphism \(f: I \to E\) such that \(d = f \circ \delta_{A}, f\) is uniquely determined by this condition for the definition of \(\delta_{A}\) gives

\[
f (x \otimes 1 - 1 \otimes x) = d x
\]

and our assertion follows from the fact that the image of \(\delta_{\mathbf{A}}\) already generates I as a left A-module (Lemma 1): hence of necessity

\[
f \left(\sum_ {i} x _ {i} \otimes y _ {i}\right) = \sum_ {i} x _ {i}. f (1 \otimes y _ {i} - y _ {i} \otimes 1) = - \sum_ {i} x _ {i}. d y _ {i}.
\]

Conversely, as the mapping \((x,y)\mapsto -x.dy\) of \(\mathbf{A}\times \mathbf{A}\) into \(\mathbf{E}\) is K-bilinear, there exists one and only one K-linear mapping \(g:\mathbf{A}\otimes_{\mathbb{K}}\mathbf{A}\rightarrow \mathbf{E}\) such that \(g(x\otimes y) = -x.dy\) ; it suffices to verify that the restriction \(f\) of \(g\) to \(\mathbf{I}\) is A-linear for the left and right A-module structures. The first assertion is obvious since \((xx') .dy = x.(x'.dy)\) ; to prove the second, note that, if \(\sum_{i}x_{i}\otimes y_{i}\in \mathbf{I}\) and \(x\in \mathbf{A}\) , then

\[
\sum_ {i} x _ {i}. d (y _ {i} x) = \sum_ {i} x _ {i}. d y _ {i}. x + \sum_ {i} (x _ {i} y _ {i}). d x
\]

but since \(\sum_{i}x_{i}y_{i} = 0\) by definition of I, this completes the proof.

We have thus defined a canonical K-module isomorphism \(f \mapsto f \circ \delta_{A}\)

\[
\operatorname{Hom} _ {(\mathbf {A}, \mathbf {A})} (\mathbf {I}, \mathbf {E}) \rightarrow \mathbf {D} _ {\mathbf {K}} (\mathbf {A}, \mathbf {E})
\]

the left hand side being the K-module of (A, A)-bimodule homomorphisms of A into E.
